\pdfoutput=1
\documentclass[letterpaper]{article} % DO NOT CHANGE THIS
\usepackage[preprint]{aaai2027}  % arXiv: show authors, no review slug
\usepackage[hyphens]{url}  % DO NOT CHANGE THIS
\usepackage{graphicx} % DO NOT CHANGE THIS
\urlstyle{rm} % DO NOT CHANGE THIS
\def\UrlFont{\rm}  % DO NOT CHANGE THIS
\usepackage{natbib}  % DO NOT CHANGE THIS AND DO NOT ADD ANY OPTIONS TO IT
\usepackage{caption} % DO NOT CHANGE THIS AND DO NOT ADD ANY OPTIONS TO IT
\frenchspacing  % DO NOT CHANGE THIS
\usepackage{algorithm}
\usepackage{algorithmic}
\usepackage{booktabs}
\usepackage{amsmath}
\usepackage{amssymb}
\usepackage{xcolor}
\usepackage{array}
\usepackage[most]{tcolorbox}

% Role colors: the same validated palette as the figures (blue=metric, orange=skill,
% green=anchor). Tables mark the proposed configuration with bold, not color.
\definecolor{mblue}{HTML}{2A78D6}
\definecolor{mbluedk}{HTML}{1C5CAB}
\definecolor{mbluebg}{HTML}{EAF2FC}
\definecolor{sorange}{HTML}{EB6834}
\definecolor{sorangedk}{HTML}{B84A1D}
\definecolor{sorangebg}{HTML}{FDF0E9}
\definecolor{agreen}{HTML}{1E7A46}
\definecolor{agreenbg}{HTML}{ECF7F0}
\definecolor{cardgray}{HTML}{5C6670}
\definecolor{cardgraybg}{HTML}{F4F5F6}
\definecolor{inkmut}{HTML}{52514E}

% Clickable cross-references and citations. The style file's hyperref guard is
% disabled in the local aaai2027.sty shipped alongside this source.
% Loaded after the palette so link colors can reference it.
\usepackage{hyperref}
\hypersetup{
  colorlinks=true,
  linkcolor=mbluedk,   % \ref to sections, figures, tables, equations, appendices
  citecolor=mbluedk,   % \cite / \citep
  urlcolor=mbluedk,
  linktoc=all
}

% Inline key-phrase highlight for Algorithm 1 (palette tint; tight box, no line-height change).
\newcommand{\algkey}[1]{{\setlength{\fboxsep}{1.5pt}\colorbox{mbluebg}{#1}}}

% Appendix card: one neutral style: light gray, dark title. Subtitle via \cardsub.
\newcommand{\cardsub}[1]{\hfill{\normalfont\footnotesize\color{inkmut}#1}}
\newtcolorbox{taskbox}[1]{enhanced, colback=cardgraybg, colframe=cardgray,
  boxrule=0.8pt, arc=2mm, left=6pt, right=6pt, top=4pt, bottom=4pt,
  fonttitle=\small\bfseries, coltitle=black, colbacktitle=cardgraybg,
  title=#1, toptitle=2.5pt, bottomtitle=0pt, before skip=6pt, after skip=6pt}

\pdfinfo{
/TemplateVersion (2027.1)
}

\setcounter{secnumdepth}{0}

\title{Who Grades the Grader?\\Co-Evolving Evaluation Metrics and Skills for Self-Improving LLM Agents}

\author{
    Xing Zhang\textsuperscript{\rm 1},
    Guanghui Wang\textsuperscript{\rm 1},
    Yanwei Cui\textsuperscript{\rm 1},
    Ziyuan Li\textsuperscript{\rm 2},
    Wei Qiu\textsuperscript{\rm 2},
    Bing Zhu\textsuperscript{\rm 2},
    Peiyang He\textsuperscript{\rm 1}\thanks{Corresponding author: peiyan@amazon.com}
}
\affiliations{
    \textsuperscript{\rm 1}AWS Generative AI Innovation Center\\
    \textsuperscript{\rm 2}HSBC Holdings Plc., HSBC Technology Center, China
}

\begin{document}

\maketitle

\begin{abstract}
Self-evolving agent systems create, revise, and retire their own skills, but every such loop assumes a reliable evaluation metric already exists. In many real applications none does. We show the metric itself can be the evolving object: our loop searches compositions of small typed drawback detectors under a full evolutionary lifecycle, selecting for agreement with a ten-item anchored reference set and regularizing by consensus over unlabeled outputs. What evolves is the function that grades one output, never the fixed task sets it is scored on, and what comes out is an inspectable expression rather than an opaque judge. It is also valid: on code generation it gains 0.21 agreement with hidden ground truth on a locked set that metric selection never reads (paired $p=0.014$), beating the bare LLM judge it contains. Validity is where safety lives: removing the anchor guards collapses the metric into a vacuous always-pass detector while removing the detector lifecycle does not, inverting the lesson from skill evolution. That collapse warns this line of work that downstream task score cannot validate a self-evolved evaluator, since the collapsed metric trains skills just as well. Task score answers only sufficiency, and an evolved metric suffices: \emph{Double Ratchet}, co-evolving the metric with a lifecycle-managed skill loop, retains 88--110\% of the lift ground truth or a hand-written rubric buys, across MBPP+, Spider~2.0-Snow, and report generation. When evolved skills gamed the report rubric, an independent judge caught it and one added detector repaired it.
\end{abstract}

\section{Introduction}

Agent self-evolution is now a standard recipe: a frozen language model is wrapped in a loop that writes skills, measures their contribution, and retires the ones that hurt \citep{zhang2026ratchet,wang2023voyager,shinn2023reflexion}. Every round consumes one expensive ingredient: an evaluation signal deciding whether an attempt passed or failed. Unit tests supply it free on code benchmarks; in deployed systems it is the bottleneck. Industry products fill the gap by hand, one rubric or demonstration at a time: Claude's managed agents have developers write outcome rubrics for a separate grader \citep{anthropic2026outcomes}, and Codex's Record \& Replay has a human demonstrate each workflow once \citep{openai2026recordreplay}. This paper automates that step.\footnote{Code: \url{https://github.com/amazon-science/Self-Evolving-Agents-Double-Ratchet}.}

The question is direct: if skills can be evolved, why not the metric? To our knowledge this is the first system to treat the evaluation metric as a first-class evolving object under a full evolutionary lifecycle in the \emph{sparse-anchor} regime: no evaluator exists to score the search, only ten anchored items with golden references. Sparse is not unsupervised, and those ten items are what the method turns on. Existing evolutionary search over objectives is scored by a downstream metric assumed valid, and learned evaluators assume exactly the ground truth we remove; here selection sees only agreement with a ten-item anchor plus consensus over unlabeled outputs, and what comes out is an \emph{inspectable expression} rather than an opaque scalar.

Two clarifications determine how every result should be read. First, \textbf{the metric is the function that scores one solver output}, and it is the only thing that evolves: not the benchmark, whose tasks and references are fixed before any loop starts, so no result here can come from an easier test. Second, \textbf{an evolved metric is judged on validity, not on the task score it produces}: agreement with a locked reference is the claim, while how well skills trained on it perform is weaker evidence, since a metric that passes everything still leaves useful practice behind. That second point is not a caveat we inherited but a result we measured, and it is strong enough to constrain how this literature should run its experiments at all.

Concretely (Figure~\ref{fig:hero}), the metric is an expression tree over a pool of small drawback detectors, each checking one failure class, synthesized from clustered failures, gated at birth, scored by marginal contribution, and retired when useless. One stance underpins the design: we rarely know what good is, but given an output we can usually find drawbacks, so a clean verdict means no known drawback was found, not certified correctness. That is what makes bootstrap possible, and mostly deterministic detectors, which fail differently from the LLM being graded, resist the shared-blind-spot collusion a bare LLM judge invites. Because a metric grading the loop that produced it is a Goodhart hazard \citep{goodhart1984problems,gao2023scaling}, anchor discipline is strict: the tiny dev set is the only supervised signal, and no part of metric evolution reads the locked set.

The hypothesis comes from the skill side: Ratchet \citep{zhang2026ratchet,zhang2026librarydrift} found unmanaged skill libraries drift and that lifecycle management is what makes skill evolution work at all. Do evaluators need the same discipline? Two ablations judged on \emph{validity} locate which ingredient carries the load, and the answer is not the one the skill literature predicts. Of our three task families, MBPP+ and Spider~2.0-Snow carry exact ground truth, so validity and sufficiency are both measurable; report generation is the intended target, having no golden metric at all, and there we add the layer a deployment would keep, an independent stronger judge comparing final outputs pairwise against the pre-evolution baseline. That audit caught a real Goodhart failure, and auditing the judge exposed a second one. Verifier-free self-improvement is feasible, but only inside an architecture that expects its own metric to be gamed and its own judge miscalibrated, and can catch both.

\begin{figure*}[t]
\centering
\includegraphics[width=\textwidth]{hero_v3.pdf}
\caption{\textbf{The evaluation metric is the evolving object; the task sets are fixed and never altered.} \emph{Center:} from a frozen solver, no metric, and ten anchored items (left), the metric loop synthesizes typed drawback detectors behind a birth gate, composes them into an expression selected by anchor agreement and consensus, and curates the pool by marginal contribution. What comes out is a readable expression, not a learned scalar. \emph{Right:} the skill loop is Ratchet \citep{zhang2026ratchet}, adopted unchanged, and the evolved metric grades its training, giving \emph{Double Ratchet}. \emph{Bottom right:} the splits metric evolution never trains on.}
\label{fig:hero}
\end{figure*}

Our contributions are as follows.
\begin{itemize}
\item \textbf{Metric evolution.} We formulate the grading function, rather than the task set, as the evolving object: anchored search over compositions of atomic drawback detectors under a full lifecycle of synthesis, birth gating, marginal-contribution scoring, and retirement.
\item \textbf{An objective that works without a downstream metric.} Recall-weighted tiny-anchor agreement paired with reliability-weighted consensus over unlabeled outputs, so consensus regularizes transfer without ever selecting a metric on its own.
\item \textbf{Where safety lives, and why task score cannot find it.} Anchor discipline, not the pool lifecycle, is load-bearing for evaluators, inverting the skill-evolution lesson; and since the collapsed metric matches the working one on task score, validating a self-evolved evaluator by downstream performance is unsound, a caution for co-evolved-judge designs generally.
\item \textbf{Double Ratchet.} Co-evolution of the metric with a skill loop, with evidence that an evolved metric suffices in place of ground truth inside the loop, plus an end-to-end Goodhart catch, its repair, and a judge audit.
\end{itemize}

\section{Related Work}

Three lines of work meet here, and in each the metric is what our method does not get for free.

\paragraph{Self-evolving agents assume the metric.}
These agents grow skill libraries, revise from failure traces, and bank reusable experience \citep{wang2023voyager,shinn2023reflexion,madaan2023selfrefine,zhao2024expel,chen2024automanual,wang2024awm,park2023generative}, a landscape two surveys map \citep{zhang2025surveyselfevolving,wang2024survey}, and a recent wave evolves and benchmarks such libraries \citep{wu2025evolver,yang2026autoskill,xia2026skillrl,alibaba2026trace2skill,zhang2026ecs,wang2026strategygenes,li2026skillsbench}. Ratchet \citep{zhang2026ratchet,zhang2026librarydrift} showed unmanaged libraries drift and that a minimal lifecycle of retirement, bounded capacity, and meta-skill guidance restores the gains; that published loop is adopted unchanged here as the skill-side component, not a contribution of this work. Others search agent code itself, in the tradition of neural architecture search \citep{hu2025automated,elsken2019nas}, or optimize prompts and pipelines \citep{khattab2024dspy,yang2023opro,yuksekgonul2024textgrad,fernando2024promptbreeder}. All optimize \emph{against} a given evaluation signal. We evolve that signal, and ask which part of the discipline carries the safety load once the evaluator moves.

\paragraph{Searched objectives are scored by a metric that exists.}
CSE-Autoloss evolves losses, AutoML-Zero whole learning algorithms, Pruner-Zero pruning metrics, AlphaEvolve program code against supplied evaluators \citep{liu2021cseautoloss,real2020automlzero,dong2024prunerzero,novikov2025alphaevolve}. The search space is analogous, primitives composed into a program, but they evolve the artifact being graded and fitness differs in kind: each is scored by a metric taken as valid, AlphaEvolve stating outright that progress must be machine-gradable, the precondition our regime lacks. This paper supplies that missing foothold: anchored agreement plus consensus in place of a downstream score. The distinction is not cosmetic. When fitness is a trusted external number, a degenerate candidate is caught by construction, since it simply scores badly; when fitness is assembled from the candidate's own agreement with a ten-item anchor and with its peers, degeneracy has to be excluded by explicit guards, which is why the guards and not the lifecycle turn out to be load-bearing here.

\paragraph{Learned graders trade the assumption for opacity.}
Reward design is delicate even in classical RL \citep{ng1999pbrs}; learned reward models are over-optimized by the policies they train, and misspecification produces systematic gaming \citep{ouyang2022rlhf,gao2023scaling,pan2022effects,skalse2022defining,amodei2016concrete}, while AI-feedback and self-rewarding setups let the model judge itself \citep{bai2022constitutional,yuan2024selfrewarding}. We share the hazard but not the structure: our grader is a composition of typed detectors optimized to predict an anchor it never trains on, so it can be read and repaired by hand, and it is never a reward the agent optimizes directly.

\paragraph{LLM judges, and auditing them.}
LLM judges are standard evaluation infrastructure and widely audited for position, verbosity, and self-preference bias \citep{zheng2023judging,dubois2024alpacafarm,gu2024surveyllmjudge,chen2024alignment}. We confine LLM judgment to three narrow guarded roles: soft labels on the ten dev items, one semantic-mismatch detector class, and a position-debiased pairwise audit outside all loops. Agent-as-judge co-evolution \citep{zhuge2024agentjudge} is closest in spirit and validates the judge through the system it grades, which our ablations show a vacuous grader also passes. Our own results add a caution to this literature that we could not have anticipated: the audit needs to know the task's output contract, or it will score a required format as a defect (\S\hyperref[sec:goodhart]{The Goodhart Episode}).

\paragraph{Reference-free evaluators for our task families.}
The evaluators nearest our tasks fix in advance what we synthesize, gate, and retire. CodeScore \citep{dong2024codescore} predicts functional correctness from the execution-labeled supervision we take to be missing, and SQLHD \citep{yang2025sqlhd} needs no ground truth but fixes a hand-authored family of metamorphic relations \citep{chen1998metamorphic}. Either could be registered as one typed op in our pool and would then compete for selection on the anchor like any other, so we read them as components rather than competitors. Agentic text-to-SQL pipelines \citep{deng2025reforce} improve the solver, an orthogonal axis: a better solver changes which failures remain, which is precisely the variable our detectability analysis isolates.

\paragraph{Benchmarks.}
MBPP+ supplies unit-tested Python tasks \citep{austin2021mbpp,liu2023evalplus}, in the execution-graded tradition of HumanEval \citep{chen2021codex}, and Spider~2.0 supplies enterprise text-to-SQL over real warehouses with execution-based grading \citep{lei2025spider2}. We use both as hard anchors for certifying an evolved metric rather than as leaderboard entries: our solver is a single LLM call with no tools or iteration, so absolute pass rates are far below what agentic systems report and are not meant to be compared with them.

\section{Problem Setup}

Let a task family come with tasks $t$ and a frozen solver mapping $t$ to an output $y$. A metric $M$ maps $(t, y)$ to a verdict in $\{\textsc{pass}, \textsc{fail}, \textsc{abstain}\}$. It is $M$ that we search over, with the tasks and references held fixed. Three disjoint splits play the usual roles with one twist each.

\begin{itemize}
\item \textbf{train}, large and unlabeled: solver outputs here expose where candidate metrics disagree or abstain, and in co-evolution the evolved metric grades these tasks for the skill loop. No reference of any kind is attached to a train task.
\item \textbf{dev}, tiny and anchored, ten items throughout: each carries a golden reference against which a teacher emits soft pass/fail labels, the only supervised signal any loop reads. Those references are not themselves a metric, grading only the ten tasks that carry them while the evolved metric must grade arbitrary outputs on tasks that have none.
\item \textbf{test}, locked against the metric: it holds the strongest available reference (unit tests, official execution comparison, the report rubric), and no part of metric evolution reads it, so metric-validity numbers are fully out of sample.
\end{itemize}

The skill loop is the one component that touches the locked set, for held-out scoring and rollback, as published \citep{zhang2026ratchet} and kept unchanged because a deployed agent rolls back against the cases it is judged on. The guarantee is thus asymmetric, and it is worth being explicit about which claim rests on which split. Metric-validity numbers are strictly out of sample: nothing in synthesis, gating, selection, or retirement ever reads the locked set, so an agreement figure there is an honest generalization estimate. Skill-side task scores are not out of sample in the same sense: they are the pass rate of a loop that rolls back on that set. They remain comparable across arms because every arm rolls back identically against the same set, which is what the head-to-head comparison uses them for, and they are never used to certify a metric. Dev is deliberately small because sparse anchoring is realistic: fresh labels in deployed systems arrive as a trickle of demos and incident reviews, not as a labeled corpus.

\subsection{Metrics as Compositions of Drawback Detectors}

An \emph{op} is an atomic drawback detector $o(t, y, c) \rightarrow \{\textsc{drawback}, \textsc{clean}, \textsc{abstain}\}$ with $c$ a shared context, checking exactly one failure class and abstaining otherwise. A task family starts from a handful covering its obvious failure modes, cheap to write by hand or have an LLM draft from the spec, and the lifecycle grows the pool from there. Ops come in three cost tiers: static ops parse the artifact, execution ops run it (sandbox for code, live warehouse for SQL), and judge ops ask an LLM one narrow question. A fixed non-evolving \emph{root} per family validates basic structure (code parses, a report has a title) and caches the parsed artifact the other ops read.

A \emph{metric expression} composes op verdicts with logical operators. Writing $v(e) \in \{1, 0, \bot\}$ for the verdict of expression $e$ on $(t, y)$, where $1$ denotes \textsc{drawback}, $0$ denotes \textsc{clean}, and $\bot$ denotes \textsc{abstain}, the grammar is
\begin{equation}
\label{eq:grammar}
\begin{aligned}
e \mathrel{::=} {}& o
  \;\;\big|\;\; \textstyle\bigvee_{i \leq k} e_i
  \;\;\big|\;\; \textstyle\bigwedge_{i \leq k} e_i
  \;\;\big|\;\; \neg e \\[2pt]
\big|\;\; & \textstyle\big[\sum_{i \leq k} v(e_i) \geq K\big]
  \;\;\big|\;\; \textstyle\big[\sum_{i \leq k} w_i\, v(e_i) \geq \theta\big],
\end{aligned}
\end{equation}
where $|$ separates the alternative node forms rather than denoting disjunction: a node is a leaf op, a disjunction (a drawback if any child finds one), a conjunction (only if all children agree), a negation, an unweighted $K$-of-$k$ vote, or a weighted vote against threshold $\theta$, so applied recursively a metric is a tree of such nodes. Abstaining children are excluded from their combinator, and a node whose children all abstain abstains itself. The output passes when the root finds no drawback, the final verdict being the fixed root conjoined with the evolved expression. Ops are pure functions of $(t, y)$, so verdicts are cached, evaluation is a set of table lookups, and a metric is reproducible from its expression string plus the registered op pool. Negation and weighted votes stay available throughout, but every metric selected in our runs uses only disjunction, conjunction, and votes.

Keeping the metric mostly deterministic is a safety choice rather than an efficiency one. A bare LLM judge shares failure modes with the LLM solver it grades, so the two can drift into agreement before any optimization pressure exists, and the agreement is invisible from the outside because there is nothing to read. A composition of typed detectors fails differently, one leaf at a time, and each leaf names the failure class it fired on. That property is what turned a rubric-gaming episode into a one-line diagnosis later in this paper.

\subsection{The Metric Loop}

Algorithm~\ref{alg:loop} sketches one round. Two signals feed synthesis: \emph{misses}, dev items the current metric passes although their soft label is fail, and \emph{gaps}, train outputs on which the op pool abstains entirely or splits. Both are clustered by failure class under a closed task-specific taxonomy, and a synthesizer authors a typed op spec for any recurring cluster. A new op must pass a birth gate: fire on at least half its cluster and stay clean on known-good outputs. Ops born from anchored misses enter \emph{active} (selectable), their birth evidence already tying them to the anchor; ops born only from unlabeled gaps enter as \emph{shadows}, which record verdicts every round but cannot be selected, since nothing yet shows they track quality rather than consensus noise. A shadow is promoted once its verdicts raise dev agreement.

Candidates come from an LLM composer that sees the op catalog, the current gaps, and an elite history of the best expressions with their shortcomings, plus mutations and crossovers of those elites. Selection keeps the best-so-far candidate by
\begin{equation}
\begin{split}
S(e) \;=\; A_{\mathrm{dev}}(e) \cdot A_{\mathrm{train}}(e)^{w} \;-\; \lambda\, C(e),
\end{split}
\label{eq:score}
\end{equation}
where $A_{\mathrm{dev}}(e)$ is $e$'s agreement with the dev soft labels, $A_{\mathrm{train}}(e)$ its agreement with the consensus verdict of the opining pool ops on unlabeled train outputs, $C(e)$ expression size, and $w, \lambda$ small constants. Consensus is the regularizer: a metric memorizing ten dev items but behaving erratically on the broad distribution loses, while a degenerate always-fail metric loses on dev. Two guards apply: selection fails closed (a candidate with no usable dev opinion is unselectable), and a validity gate drops candidates that pass, fail, or abstain on everything. Per-op fitness is the leave-one-out marginal of the incumbent expression, and non-positive ops retire after a grace period. Every constant, including the exponent $w$, the complexity charges, and the grace length, is listed in \hyperref[app:lang]{Appendix~C}.

Both agreement terms are weighted, each weighting a safety choice rather than a tuning knob. \emph{Dev agreement is recall-weighted}: $A_{\mathrm{dev}}$ is the weighted mean of the two per-class recalls, $(w_{\textsc{f}} r_{\textsc{fail}} + w_{\textsc{p}} r_{\textsc{pass}})/(w_{\textsc{f}} + w_{\textsc{p}})$ with $w_{\textsc{f}} = 2$, $w_{\textsc{p}} = 1$, over the items the metric opines on, so a missed drawback costs twice a false alarm and no candidate reaches $A_{\mathrm{dev}} = 1$ by acing the easier class alone. \emph{Consensus is reliability-weighted}: each opining op's vote is scaled by $1 + \kappa \max(0, m_o)$ for $m_o$ that op's anchored leave-one-out marginal, so anchor-proven ops lead and unproven ops vote near-neutrally. Otherwise a bloc of correlated detectors, cheap to synthesize and liable to share a blind spot, would define the very pseudo-label the regularizer rewards agreeing with. $A_{\mathrm{train}}$ is thus a transfer regularizer pulled toward the anchor, not a popularity vote, and entering multiplicatively against a fail-closed $A_{\mathrm{dev}}$ it never selects a metric alone.

\begin{algorithm}[t]
\caption{One round of metric evolution. Only the audit step reads the locked test set, and it never feeds back. Constants and the full lifecycle are in \hyperref[app:lang]{Appendix~C}.}
\label{alg:loop}
\begin{algorithmic}[1]
\STATE {\color{mbluedk}\textbf{Sense}} \ {\footnotesize\color{inkmut}\emph{(where is the metric wrong?)}}
\STATE \quad misses $\gets$ dev items the incumbent passes but the soft label fails
\STATE \quad gaps $\gets$ train outputs where the pool abstains or splits
\STATE {\color{mbluedk}\textbf{Grow}} \ {\footnotesize\color{inkmut}\emph{(expand the op pool)}}
\STATE \quad cluster misses $\cup$ gaps by failure class; author one typed op per recurring cluster
\STATE \quad birth gate: fires on $\geq$ half its cluster and clean on known-good $\Rightarrow$ admit (\algkey{active if anchored, else shadow})
\STATE {\color{mbluedk}\textbf{Select}} \ {\footnotesize\color{inkmut}\emph{(recombine, keep the best)}}
\STATE \quad candidates $\gets$ LLM compositions over the pool, plus \algkey{mutations and crossovers of elites}
\STATE \quad keep $\arg\max_e S(e)$ over candidates and the incumbent, behind the validity gate (Eq.~\ref{eq:score})
\STATE {\color{mbluedk}\textbf{Curate}} \ {\footnotesize\color{inkmut}\emph{(merit-based lifecycle)}}
\STATE \quad retire ops with non-positive \algkey{leave-one-out marginal}; promote shadow ops that raise dev agreement
\STATE {\color{agreen}\textbf{Audit}} \ {\footnotesize\color{inkmut}\emph{(measurement only, never a training signal)}}
\STATE \quad report agreement against the locked test set
\end{algorithmic}
\end{algorithm}

\subsection{Co-Evolution with the Skill Loop: Double Ratchet}

The skill loop adopted here \citep{zhang2026ratchet} needs one thing our setting does not supply: a grader for its training tasks. Co-evolution hands it the current best evolved metric, giving \emph{Double Ratchet} (Figure~\ref{fig:hero}). The loops alternate under a fixed front-loaded curriculum, metric phases longest early when coverage is worst: the evolved metric grades train attempts, and failures become capsules whose error text feeds skill synthesis. Held-out evaluation and the rollback anchor stay pinned to the locked test set, which the evolved metric never touches, so a corrupted metric slows skill learning without corrupting the measurement. The comparison is the reference skill loop, identical machinery handed the ground-truth metric (or best available rubric) for free.

\paragraph{Two ablations, one per safety ingredient.}
The \textbf{naive} arm disables the anchor guards (fail-closed anchoring, the validity gate, and skill-loop rollback) with the detector lifecycle fully on, the collusion-prone configuration unanchored agent-judge co-evolution implicitly runs. The \textbf{no-lifecycle} arm does the opposite, keeping the guards but disabling the birth gate and merit-based retirement so every synthesized detector enters active and nothing retires: the metric-side analog of an unmanaged skill library. The two arms are deliberately not a factorial. Each isolates one ingredient against an otherwise complete system, which is what a claim about which ingredient is load-bearing needs; a fully crossed design would add an arm with neither guard nor lifecycle whose collapse is already implied by the naive arm.

\paragraph{Independent final judge.}
Where no accurate metric exists, rubric scores are references and not truth, so a final audit sits outside all loops: a stronger LLM compares each final output pairwise against the pre-evolution baseline, judged twice with positions swapped, a win counting only when both orders agree. It mimics the human acceptance step a deployment keeps, and exposed a proxy-gaming failure the loops could not see. The judge is an audit and never a training signal: no gradient, capsule, hint, or selection decision anywhere in the system reads it.

\section{Experimental Setup}

\paragraph{Task families.}
\textbf{MBPP+} \citep{austin2021mbpp,liu2023evalplus}, Python function synthesis with hidden unit tests as the hard anchor, and \textbf{Spider~2.0-Snow} \citep{lei2025spider2}, enterprise text-to-SQL over real Snowflake warehouses anchored by official execution-result comparison, are the validation tasks: their exact ground truth makes both metric validity and the reference-loop match measurable. \textbf{Report generation} is the target task, a deployment-style workload in which a pipeline writes analyst-report sections from audited evidence cards. No golden metric exists for it: the available rubric score (RAQS) plus golden demonstrations is a partial reference signal, not ground truth. Task cards are in \hyperref[app:tasks]{Appendix~A}; tables abbreviate the tasks as MBPP+, Spider, and Report.

\paragraph{Hard-subset construction.}
For MBPP+ and Spider~2.0 we screen every task with five frozen-solver samples, drop those passed on every sample (no lift room), and stratify the rest by screen pass rate. Splits are 60/10/40 for MBPP+, 59/10/40 for Spider~2.0, and 73/10/48 for reports, whose tasks derive from golden demonstrations and need no screening. Dev always holds ten items, five clear failures and five clean positives. The solver is a single direct LLM call with no tools, so absolute numbers are not comparable to leaderboard agents.

\paragraph{Roles and models.}
All experiments use Claude models: every in-loop role (solver, skill synthesizers, teachers, judge ops) is Claude Opus~4.7, the sole exception being the report task's final judge, the stronger Claude Opus~4.8 kept outside all loops. A same-family teacher is not a circular self-grader: it compares candidate outputs against golden references on dev rather than judging unaided, and never sees the locked anchor. Skill loops run 100 rounds; co-evolution interleaves metric phases of 15/8/5/2 rounds between four 25-round skill phases, matching the reference loop's skill budget. Every role is bound per call from a role-to-model table, so repeating any experiment under a different or mixed model assignment is a configuration change rather than a code change; we hold that table uniform here to keep the comparisons clean (\hyperref[app:repro]{Appendix~H}).

\section{Results}

The findings run along two axes, in a deliberate order: first \emph{metric validity}, agreement with a locked reference no part of metric evolution reads, which is the claim and which the guard ablations then localize; then \emph{task score}, which we show cannot certify a metric and so use only to ask whether an evolved one suffices in place of ground truth. Each run uses three seeds, reported as mean and sample standard deviation, with loop comparisons adding the seed-level difference and a 95\% bootstrap interval. Three seeds and held-out sets of tens of items make those means coarse, so metric claims are also tested \emph{paired}, exact McNemar over matched per-item verdicts, since evolution rescores the same locked items and two rounds differ only in the grader. That test pools three seeds over one 40-item set, so its unit is the seed-item verdict, not an independent task: it shows the change of grader is not noise, not a task-level population effect. Per-round noise runs near five points, so we report peak and final-rounds values (mean of the last ten) against each run's round-0 baseline. Peak is a locked-set argmax and so optimistic; the final-rounds column is fixed in advance and is the unselected estimate, which is why the paired test is run at both. Bold rows are our proposed configuration.

\subsection{The Evolved Metric Lifts and Transfers}

Figure~\ref{fig:results}c is the mechanism made concrete: every distinct metric the MBPP+ loop selected, in order, as the expression it actually is. The objective rises from 0.16 to 0.52 while the expression \emph{shrinks} from six ops to three, the lifecycle working as designed rather than a search accumulating checks, and the birth gate is the busy part, 998 candidate detectors rejected against 2 admitted on this seed. Table~\ref{tab:oplife} puts that ratio in context across tasks: the pools stay small, the gate does most of the work, and the selected expression is one to three leaves everywhere. What rises monotonically is what selection optimizes: best-so-far selection makes $S(e)$ never fall, on all nine runs (full 100-round budget, early stopping off, last improvement by round 63). Held-out agreement, which selection never sees, is not monotone: five of nine runs end below their own mid-run peak, the ordinary generalization gap of optimizing a ten-item proxy, and why we report the locked column instead of selecting on it. Convergence is always to a composition: \texttt{(any spec\_mismatch crash returns\_not\_print\_only)} on MBPP+ and two-op forms such as \texttt{(any missing\_group\_by spec\_sql\_mismatch)} elsewhere (\hyperref[app:artifacts]{Appendix~E}).

\begin{table}[t]
\centering
\footnotesize
\setlength{\tabcolsep}{2.5pt}
\begin{tabular}{lccccc}
\toprule
Task & Seed & Born & Rejected & Pool & In metric \\
\midrule
MBPP+ & 10 & 2--4 & 997--1014 & 12--14 & 3 \\
Spider & 9 & 24--30 & 612--785 & 33--39 & 1--2 \\
Report & 11 & 11--15 & 172--188 & 22--26 & 2 \\
\bottomrule
\end{tabular}
\caption{\textbf{The birth gate rejects orders of magnitude more detectors than it admits, and selection composes only one to three survivors.} Metric-loop op lifecycle, anchored arm, 100 rounds. \emph{Seed}: hand-authored detectors. \emph{Born}: synthesized and admitted during the run. \emph{Pool}: seed plus born. \emph{In metric}: leaves in the final expression. Three-seed min--max; a bare number means all seeds agree. Skill-side counterpart: Table~\ref{tab:skilllife} in \hyperref[app:lifecycle]{Appendix~F}.}
\label{tab:oplife}
\end{table}

\paragraph{What the lift is measured against.}
These compositions beat the bare LLM judge op they contain, the natural baseline for a grader built without ground truth: peak agreement $0.625\pm0.050$ versus $0.55\pm0.04$ on MBPP+ and $0.500\pm0.026$ versus $0.45\pm0.07$ on Spider (no ground truth exists to score a standalone judge on Report). Round~0 brackets the evolved metric from the other side, being the hand-authored seed composition, at $0.417\pm0.058$ on MBPP+, so the $+0.21$ over it is lift over static engineering rather than lift over nothing. Two baselines would split the remaining credit and neither is run here: a prompt-tuned judge on the same ten items, which would separate anchored composition from anchored prompting, and random composition in place of the LLM composer, which would separate the composer from the search space it draws on.

\paragraph{Three tasks, three roles.}
The tasks are not three attempts at one claim. MBPP+ is the positive case, the only measurable held-out lift; Report starts at ceiling from round 0, where subset seeding already supplies a near-optimal composition, so it carries sufficiency and the safety ablations rather than lift; Spider sits near chance despite a robust objective climb, the detectability boundary taken up in \S\hyperref[sec:detect]{The Detectability Spectrum}. Absolute agreement is not comparable across the three, since each sits on its own locked set with its own reference and class balance, which is why the within-task Spider contrast carries the detectability claim and a cross-task ranking would not.

\paragraph{The lift survives a paired test.}
On MBPP+ the paired test puts the lift beyond seed noise: across 120 matched verdicts, evolution corrects 61 items against 36 broken at peak ($p=0.014$) and 36 against 20 at the end ($p=0.044$), every seed net positive. It finds nothing on Spider, that task's regime rather than a sample-size artifact, since the same test on the same set size detects the MBPP+ effect. Nor is the ten-item anchor fragile: replaying report selection offline on 200 stratified subsamples per size, rescoring every candidate the live run considered, gives held-out agreement $0.871\pm0.063$, $0.882\pm0.051$, and $0.874\pm0.041$ at four, six, and eight anchored items against 0.854 at the full ten. Smaller anchors are noisier but their means do not fall, so recall weighting plus the consensus regularizer keep even a four-item anchor off degenerate metrics. The standard deviations tell the operational story: any single four-item anchor is a lottery, while the mean over draws shows selection is not leaning on one marginal item (protocol in \hyperref[app:negative]{Appendix~G}).

\begin{figure*}[t]
\centering
\includegraphics[width=\textwidth]{results_v3.pdf}
\caption{\textbf{An evolved metric buys most of the lift ground truth would have, and what evolves stays readable.} \textbf{(a)} Lift Double Ratchet retains as a fraction of the reference loop's; 100\% means the evolved metric bought everything ground truth would have. \textbf{(b)} Locked-set learning curves (mean $\pm$ std, three seeds): reference loop orange, Double Ratchet blue, frozen solver dotted. \textbf{(c)} Every distinct metric the MBPP+ loop selected, in order. The objective rises while the expression shortens from six ops to three; the locked column is reported, never read. Numbers from Tables~\ref{tab:main} and \ref{tab:perloop}.}
\label{fig:results}
\end{figure*}

\subsection{Which Guard Carries the Load, and Why Task Score Cannot Tell}
\label{sec:ablation}

Table~\ref{tab:ablation} contrasts three arms of the report metric loop, where the anchor is softest and collusion pressure highest. Disabling the anchor guards (naive) collapses the metric on three of three seeds: selection latches onto a detector that almost never fires, fail-open scoring reports a vacuous perfect objective, and the metric passes 0.97--1.00 of everything it grades. Disabling the lifecycle does not: objectives and held-out agreement stay in the anchored band, and although the pool grows (23--31 ops against 22--26 anchored), the extra ops are inert, since an op never selected never grades anything. This inverts the skill-evolution literature's emphasis: for skills, lifecycle management is the key finding; for evaluators under anchored selection, anchor discipline carries the safety load, so the metric-side analog of library drift is anchor drift, not pool drift. Structurally, a junk skill gets \emph{routed} into prompts and hurts, while a junk op matters only if selection picks it. The lifecycle is not therefore worthless: it keeps the pool and its per-round verdict cost bounded (Table~\ref{tab:oplife}). It is simply not what stands between this system and a worthless grader.

\begin{table}[t]
\centering
\footnotesize
\setlength{\tabcolsep}{2.2pt}
\begin{tabular}{lcccc}
\toprule
Arm & Objective & Held-out & Train & Outcome \\
\midrule
\textbf{anchored} & \textbf{0.865$\pm$0.002} & \textbf{0.830$\pm$0.012} & \textbf{.78--.82} & \textbf{composes} \\
naive & 1.000 (vacuous) & fail-open & .97--1.0 & collapses \\
no-lifecycle & 0.896$\pm$0.072 & 0.868$\pm$0.061 & n/a & no collapse \\
\bottomrule
\end{tabular}
\caption{\textbf{Removing the anchor guards collapses the metric into a vacuous always-pass grader; removing the lifecycle does not.} Metric-validity ablations on report generation, final-rounds values over three seeds. \emph{Objective}: the selection score of Eq.~\ref{eq:score}. \emph{Held-out}: agreement with the locked reference. \emph{Train}: fraction of training attempts the metric passes. The naive arm scores 1.000 because fail-open scoring credits a detector that never fires; it passes essentially everything, so its held-out entry is a behavior rather than a number. Which guard each arm disables is specified in \hyperref[app:lang]{Appendix~C}.}
\label{tab:ablation}
\end{table}

\begin{figure*}[t]
\centering
\includegraphics[width=\textwidth]{dissoc_v1.pdf}
\caption{\textbf{Metric validity and downstream task score dissociate, so task score cannot certify a self-evolved grader.} The same three report-generation arms, measured two ways. \textbf{(a)} Metric validity on an absolute 0-to-1 scale against a locked reference no metric loop reads: the arms separate sharply, and the naive arm is vacuous rather than merely worse. \textbf{(b)} The task score of the skill loop each grader trained, on a zoomed axis: all three land within .038, and the worthless grader is the higher of the two co-evolved arms. Both readings are correct, which is the problem: certifying a grader by the performance of the policy it trained would have certified the hatched bar. Values from Tables~\ref{tab:ablation} and \ref{tab:main}.}
\label{fig:dissoc}
\end{figure*}

The collapsed arm also carries a warning past this paper, drawn out in Figure~\ref{fig:dissoc}. Its metric is worthless, yet its co-loop scores as well as the anchored one and on two tasks slightly better (Table~\ref{tab:main}). Neither number is wrong: an always-pass metric reduces training to unfiltered practice, which on its own improves skills, and measurement never passes through the metric, so the score survives what the grader does not. The consequence is methodological. \textbf{Downstream task performance cannot validate a self-evolved evaluator}, and the experiment this literature reaches for first, co-evolve a judge with a policy and report that the policy improved, would have certified a grader that passes everything. Validating the evaluator needs a held-out reference it never reads, as above, which is also what a deployment triaging or gating on it depends on.

The asymmetry is worth stating in both directions, because it is easy to over-read. A high task score does not certify the grader, as the naive arm shows. But a low task score does not condemn one either, since a grader can be right and still drive a loop into a capability ceiling it cannot reach, which is what happens on Spider below. Task score and metric validity are simply different measurements of different objects, and the practice of substituting the first for the second is what we are arguing against.

\subsection{Co-Evolution: Is the Evolved Metric Sufficient?}

\begin{table}[t]
\centering
\footnotesize
\setlength{\tabcolsep}{3pt}
\begin{tabular}{lccc}
\toprule
 & MBPP+ & Spider & Report \\
\midrule
Peak: skill & 0.700$\pm$0.025 & 0.483$\pm$0.038 & 0.850$\pm$0.010 \\
Peak: \textbf{co} & \textbf{0.717$\pm$0.038} & \textbf{0.458$\pm$0.038} & \textbf{0.812$\pm$0.006} \\
Peak: co naive$^\dagger$ & 0.742$\pm$0.014 & 0.458$\pm$0.029 & 0.841$\pm$0.003 \\
\midrule
Lift retention (co) & 106\% & 110\% & 88\% \\
$\Delta$peak (co $-$ skill) & $+.02$ & $-.03$ & $-.04$ \\
\quad 95\% CI & $[-.03,.06]$ & $[-.08,.03]$ & $[-.05,-.03]$ \\
\midrule
Improved: skill & 16/23 & 4/12 & 99\% \\
Improved: \textbf{co} & \textbf{19/23} & \textbf{6/12} & \textbf{99\%} \\
\bottomrule
\end{tabular}
\caption{\textbf{Double Ratchet retains 88--110\% of the lift ground truth or a hand-written rubric would have bought.} Held-out task score, three seeds, 40--48-item locked sets. \emph{skill}: the reference loop, trained on ground truth (MBPP+, Spider) or the best-available rubric (Report). \emph{co}: Double Ratchet. \emph{co naive}: the anchor-guards-off ablation, seed-matched to \emph{co}, whose metric is vacuous ($^\dagger$; Table~\ref{tab:ablation}) yet scores as well, which is why this table answers sufficiency but cannot certify a grader (Figure~\ref{fig:dissoc}). \emph{Lift retention}: co's peak lift over round~0 as a fraction of skill's. \emph{$\Delta$peak}: the seed-level difference with a 95\% bootstrap CI, containing zero on both ground-truth tasks. \emph{Improved}: never-passed tasks solved at least once, or (Report) sections ending above round~0.}
\label{tab:main}
\end{table}

\begin{table}[t]
\centering
\footnotesize
\setlength{\tabcolsep}{3pt}
\begin{tabular}{llccc}
\toprule
Task & Loop & Init & Peak & End \\
\midrule
MBPP+ & skill & 0.275$\pm$0.050 & 0.700$\pm$0.025 & 0.592$\pm$0.039 \\
 & \textbf{co} & \textbf{0.267$\pm$0.038} & \textbf{0.717$\pm$0.038} & \textbf{0.652$\pm$0.034} \\
 & co naive & 0.317$\pm$0.058 & 0.742$\pm$0.014 & 0.632$\pm$0.023 \\
\midrule
Spider & skill & 0.317$\pm$0.038 & 0.483$\pm$0.038 & 0.327$\pm$0.010 \\
 & \textbf{co} & \textbf{0.275$\pm$0.043} & \textbf{0.458$\pm$0.038} & \textbf{0.323$\pm$0.038} \\
 & co naive & 0.292$\pm$0.052 & 0.458$\pm$0.029 & 0.343$\pm$0.028 \\
\midrule
Report & skill & 0.562$\pm$0.001 & 0.850$\pm$0.010 & 0.802$\pm$0.024 \\
 & \textbf{co} & \textbf{0.557$\pm$0.006} & \textbf{0.812$\pm$0.006} & \textbf{0.767$\pm$0.015} \\
 & co naive & 0.559$\pm$0.003 & 0.841$\pm$0.003 & 0.811$\pm$0.015 \\
\bottomrule
\end{tabular}
\caption{\textbf{The arms stay ordered at the unselected End column exactly as at peak, so retention is not an artifact of the peak argmax.} Per-loop scores behind Table~\ref{tab:main}, mean $\pm$ std over three seeds. \emph{Init} is round~0, \emph{Peak} the best round on the locked set, \emph{End} the mean of each run's last ten rounds. Frozen single-call baselines are $\approx$0.27 (MBPP+), 0.285 (Spider~2.0), and 0.56 (Report). Bold rows are our proposed configuration.}
\label{tab:perloop}
\end{table}

With validity established independently (\S\hyperref[sec:ablation]{Which Guard Carries the Load}), the remaining question is sufficiency: can a metric evolved from ten reference examples drive a skill loop that ground truth would otherwise drive? It can. Double Ratchet retains 106\%, 110\%, and 88\% of the reference loop's lift on MBPP+, Spider, and Report (Table~\ref{tab:main}; Figure~\ref{fig:results}a--b): within noise on the ground-truth tasks, and a small gap where the reference is itself a \emph{partial rubric}, which is why the independent final judge below carries more weight there. Table~\ref{tab:perloop} gives the underlying init / peak / final-rounds triples, and the ordering of arms at the unselected final-rounds column matches the ordering at peak, so retention is not an artifact of taking an argmax on the locked set.

Spider's absolute lift is smallest for \emph{both} loops in Figure~\ref{fig:results}b, one ceiling binding them alike: a single-call solver's residual failures mostly need live exploration no guidance-text skill expresses. Aggregates also hide a task-level result: both loops solve held-out problems the frozen solver \emph{never} solves, and on both hard-anchor tasks Double Ratchet converts \emph{more} of them than the reference loop (Table~\ref{tab:main}, \emph{Improved}). That settles how good a metric must be to \emph{suffice} here: not very. Spider's metric agrees with ground truth only at $0.500\pm0.026$ yet its co-loop retains the full reference lift, because the metric's training role is \emph{directional}: it decides which attempts become failure capsules, whose concrete error text does not depend on the label being perfect. Metric-graded training accuracy stays in a realistic band (0.78--0.82 on Report, 0.41--0.83 on Spider), so the metric is really grading, and a weak one slows learning without corrupting measurement.

That robustness has a sharp edge, visible on Spider at the seed level: two of three co-evolution seeds selected a weak single-op metric with train agreement 0.41--0.44, yet reached peak evaluations of 0.45--0.50, indistinguishable from the sibling seed carrying a two-op metric (\hyperref[app:negative]{Appendix~G}). A skill loop is a forgiving consumer of grades. The reasons to build a metric well therefore lie elsewhere, in the uses a skill loop does not exercise: triage, where a grader decides which outputs a human sees; gating, where it decides which ship; and repairability, where its structure decides whether a failure can be localized at all. The next subsection puts repairability under load.

\subsection{The Goodhart Episode and the Final-Judge 2$\times$2}
\label{sec:goodhart}

The report task produced the paper's most instructive sequence (Table~\ref{tab:judge}), the clearest evidence that an inspectable metric is worth having. The rubric's metric-discipline dimension counts inline evidence tags, and the first skill runs raised the rubric score by $+0.26$ partly by gaming exactly this: evolved skills wrote the tag \emph{in place of} the number (about 30\% of tags at peak rounds had no value beside them) and invented confident forecasts for a style dimension. The independent judge preferred the baseline in 88\% of decided pairs, citing unrendered placeholders and fabrication. This is a fact about the setup rather than about our system: the rubric was the best available reference signal on this task, hand-written by the pipeline's authors, and gamed within a hundred rounds by skills optimizing against it. Any design that treats such a rubric as ground truth inherits that outcome silently.

The repair was one coupled change: a vocabulary-aware value-erasure check in the capsule gate (a tag whose registered metric has a numeric value but no numeral nearby is a defect) plus the failure hints that teach value-plus-tag, the hint being how the detector's verdict reaches synthesis. We claim the pair, not the detector alone, since separating them needs a detector-only and hints-only rerun we did not do. The rerun cut erased tags to about 1\% while the rubric score \emph{rose} to $0.850\pm0.010$, so nothing was traded away for the fix. Reading the metric is what localized the defect to a nameable check, an intervention a learned scalar does not admit: a reward model that had absorbed the same rubric would have shown a number going up, with no place to put the repair.

The judge itself then needed auditing. The pipeline requires raw evidence tags in its output, a format contract the original generic rubric did not know, so its loss rationales blamed that syntax and the repair barely moved the win rate (0.122 to 0.126). A task-aware rubric stating the contract, still treating valueless tags and fabricated figures as defects, re-judged the same stored pairs: down that column the fix moves the win rate from 0.515 (271 of 526 decided) to 0.770 (435 of 565), real content quality, while across each row the jump measures convention-blindness. That rubric was written after reading the generic judge's rationales, so it diagnoses the audit rather than independently confirming the repair; its pre-repair cell is the control, both cells re-judging identical stored outputs. Both directions matter: self-evolved metrics need an outer audit, and the audit needs the task contract. A deployment that had stopped at the generic column would have concluded the repair did nothing and reverted it.

\begin{table}[t]
\centering
\small
\setlength{\tabcolsep}{4.5pt}
\begin{tabular}{lcc}
\toprule
 & Generic judge & Task-aware judge \\
\midrule
pre-repair (gamed proxy) & 0.122 & 0.515 \\
post-repair (erasure fixed) & 0.126 & \textbf{0.770} \\
\bottomrule
\end{tabular}
\caption{\textbf{Reading down a column shows what the repair bought; reading across a row shows how much the judge penalized a required output format.} Final-judge win rate of evolved report outputs over their pre-evolution baselines, over decided pairs (both position orders agreeing); ties are 39--50\% of the 864 pre-repair and 1{,}056 post-repair pairs. All four cells re-judge the same stored outputs. Both rubrics and the pairing protocol are in \hyperref[app:judge]{Appendix~D}.}
\label{tab:judge}
\end{table}

\subsection{The Detectability Spectrum}
\label{sec:detect}

One variable organizes the results above: how mechanically detectable the solver's failures are. MBPP+ failures are crashes and wrong outputs cheap ops see directly, and reports have a closed failure taxonomy behind their anchor (\hyperref[app:setup]{Appendix~B}), so transfer is strong on both. Spider~2.0 turns that into an in-task controlled contrast. Under our first prompt most failures were compile errors and the evolved metric reached 0.85 held-out agreement, but grounding the prompt with typed schemas raised the solver's baseline by half and left \emph{residual} failures that are semantically wrong values under clean execution, invisible to deterministic ops, so agreement fell to $0.500\pm0.026$ even as the objective lifted. Same method, same task, different failure distribution: a semantic regime shifts the burden onto judge ops and outer audits, and the sharper the solver, the more of this regime is left. What evolves stays legible either way, which is what made the repair possible.

Figure~\ref{fig:detect} lays the three task families on that axis and keeps the controlled comparison separate from the cross-task placement, because only the former is a measurement. The placement of MBPP+, Report, and Spider on the spectrum is an interpretation of where each family's failures sit; their agreement numbers are not comparable to one another, since each is computed against a different locked reference on a different set. The Spider band underneath is the actual evidence: one task, one method, two solver prompts, and the metric's held-out agreement moving from 0.85 to 0.500 as the residual failures move from mechanical to semantic.

\begin{figure*}[t]
\centering
\includegraphics[width=\textwidth]{detect_v1.pdf}
\caption{\textbf{Whether an evolved metric transfers is governed by how mechanically detectable the solver's failures are.} \emph{Top:} the three task families on the axis from mechanically checkable failures to purely semantic ones, each with its evolved metric's peak held-out agreement. Those three numbers sit on three different references, so the placement is an interpretation, not a measurement. \emph{Middle, the measurement:} a within-task Spider~2.0-Snow contrast. Grounding the solver prompt with typed schemas pushes the residual failures from compile errors into clean-executing wrong values, and the same metric loop's held-out agreement falls from 0.85 to 0.500. \emph{Bottom:} what each end of the axis implies for practice.}
\label{fig:detect}
\end{figure*}

The practical reading is a scoping rule. Where failures are mechanically checkable, cheap deterministic detectors carry the metric and an evolved composition transfers; where they are purely semantic, the composition still evolves and still names what it finds, but its agreement with truth is bounded by the judge ops inside it, and the outer audit stops being optional. The uncomfortable corollary is that this boundary moves in the direction of harder, not easier: every improvement to the solver retires the failures a static op could catch and leaves behind the ones only judgment can.

\section{Discussion}

\paragraph{The anchor cannot be manufactured.}
Evolution expands coverage but never creates ground truth. With no anchor the loop optimizes noise, which is what the naive arm demonstrates from the inside: remove the guards that tie selection to ten real labels and the search finds the cheapest way to look perfect. Anchor \emph{quality} therefore matters as much as its existence, and hardening soft anchors into checkable ones is the highest-leverage investment available. Ten items is few; ten items with golden references and a strict soft-labeling protocol were enough on all three families, and the subsample replay says four would often have been.

\paragraph{Validity and sufficiency are different questions.}
The two axes of the results section are not two views of one number. Sufficiency asks whether a loop trained on the evolved metric reaches where ground truth would take it, and the answer is largely yes, because a skill loop consumes grades directionally. Validity asks whether the grader is right, and it is the question every use outside a training loop turns on. Conflating them is not a philosophical error but a measurable one: Figure~\ref{fig:dissoc} shows a grader that answers the first well and the second not at all.

\paragraph{Where the method helps.}
Metric evolution buys the most where the solver's failures are mechanically detectable, as the within-task Spider contrast quantifies (\S\hyperref[sec:detect]{The Detectability Spectrum}). It buys least where they are semantic, and there its value shifts from agreement to legibility: an expression that names its verdicts supports triage and repair even when it barely beats chance at predicting truth.

\section{Limitations}

This is a mechanism study, not a scaling result. Three task families spanning the detectability axis including its negative end, one model family, rounds in the low hundreds, and held-out sets of tens of items, though graded every round by every candidate (52{,}600 item-level verdicts against the locked reference and 5{,}743 birth-gate decisions). Nothing in the design is model-specific, and every LLM role is bound from a configuration table, but we have not run a cross-family replication and cannot claim one.

Several specific gaps are worth naming. Teacher quality is held fixed, so how anchor quality degrades with a weaker teacher is untested. The paired McNemar test pools seeds over a single 40-item locked set, so it licenses a claim about graders, not about tasks. Peak values are locked-set argmaxes and therefore optimistic, which is why the final-rounds column, fixed in advance, is reported beside them everywhere. The Goodhart repair is claimed as a coupled pair of detector and hints, since we did not run the detector-only and hints-only arms that would separate them. The task-aware rubric was written after reading the generic judge's rationales, so it is a diagnosis and not an independent confirmation. Guidance-text skills cannot substitute for capabilities that need interaction, such as live schema exploration, which is the ceiling Spider's absolute lift runs into for both loops. And metamorphic detectors \citep{chen1998metamorphic}, the natural next op family for exactly the semantic regime where our current pool is weakest, are built but unexercised here.

\section{Conclusion}

The evaluation metric can itself be the evolving object, and it evolves into something valid: $+0.21$ held-out agreement on MBPP+ against a locked reference no part of selection reads, ahead of the bare judge it contains, and beyond seed noise under a paired test. Anchor discipline, not the pool lifecycle, makes that hold, inverting the lesson the skill-evolution literature would predict, and the arm without it warns that task score certifies nothing, a worthless metric training skills just as well. Read for sufficiency rather than validity, Double Ratchet retains 88--110\% of the supervised lift across code, enterprise SQL, and report writing. On the one task where no golden metric exists, an outer judge caught the loop gaming its own reference, one detector plus its hints repaired it without cost to the score, and auditing the judge caught the audit mistaking a required output format for a defect. Who grades the grader: an anchor it predicts but never sees, and an outside judge that must itself be checked.

\documentclass{article}


\usepackage{PRIMEarxiv}

\usepackage[numbers, sort, super]{natbib}
\usepackage{natmove}
\usepackage[utf8]{inputenc} % allow utf-8 input
\usepackage[T1]{fontenc}    % use 8-bit T1 fonts
\usepackage[hidelinks]{hyperref} 
\usepackage{url}            % simple URL typesetting
\usepackage{booktabs}       % professional-quality tables
\usepackage{amsfonts}       % blackboard math symbols
\usepackage{nicefrac}       % compact symbols for 1/2, etc.
\usepackage{microtype}      % microtypography
\usepackage{fancyhdr}       % header
\usepackage{graphicx}       % graphics
\usepackage{algorithm}      % pseudocode
\usepackage[noend]{algpseudocode}  % pseudocode, and noend means don't put 'end procedure'
\usepackage{amsmath}        % \gets
\usepackage{setspace}       % increasing algorithm vertical spacing
\graphicspath{{media/}}     % organize your images and other figures under media/ folder
\usepackage{dsfont}         % indicator variable
\usepackage{subcaption}     % subtables
\usepackage[table]{xcolor}  % table background color
\usepackage{bibunits}
\usepackage{everypage, eso-pic}      % JL

% Used for pseudocode, SEE: https://tex.stackexchange.com/a/163779
\makeatletter
\def\BState{\State\hskip-\ALG@thistlm}
\makeatother

% Used for subsubsection, SEE: https://tex.stackexchange.com/a/667422
\renewcommand{\sectionautorefname}{Section}
\renewcommand{\subsectionautorefname}{Section}
\renewcommand{\subsubsectionautorefname}{Section}

\renewcommand*\thefootnote{\alph{footnote}}
\renewcommand{\labelitemi}{$\bullet$} 
\renewcommand{\labelitemii}{-}        
\renewcommand{\labelitemiii}{$\circ$}
%Header
\pagestyle{fancy}
\thispagestyle{empty}
\rhead{ \textit{ }}


% Update your Headers here
%\fancyhead[LO]{Language agents for scientific knowledge}
% \fancyhead[RE]{Firstauthor and Secondauthor} % Firstauthor et al. if more than 2 - must use \documentclass[twoside]{article}

%% JL fancy backgrounds below

% Define backgrounds for pages
\newcommand\FirstPageBackground{%
  \AddToShipoutPictureBG*{%
    \put(0,0){%
      \parbox[b][\paperheight]{\paperwidth}{%
        \vfill
        \centering
        \includegraphics[width=\paperwidth,height=\paperheight,%
                         keepaspectratio]{Figures/wikicrow_background-title.pdf}%
        \vfill
      }}}}

\newcommand\OtherPagesBackground{%
  \AddToShipoutPictureBG*{%
    \put(0,0){%
      \parbox[b][\paperheight]{\paperwidth}{%
        \vfill
        \centering
        \includegraphics[width=\paperwidth,height=\paperheight,%
                         keepaspectratio]{Figures/wikicrow_background-main.pdf}%
        \vfill
      }}}}

% Add background to first page
\AtBeginShipoutFirst{\FirstPageBackground}

% Add background to other pages starting from the second page
\AddEverypageHook{%
  \ifnum\value{page}>0 %
    \OtherPagesBackground
  \fi}

%%


% Update footers here
%#\fancyfoot[C]{\textbf{CONFIDENTIAL}}

%\fancypagestyle{firstpage}{%
%  \fancyhead{}
%  \renewcommand\headrule{}
%  \fancyfoot[C]{\textbf{CONFIDENTIAL}}
%}

%% JL page numbers
\fancyfoot[R]{%
  \makebox[0pt][l]{%
    \hspace*{-0.225in}%
    \raisebox{-0.135in}{%
    \makebox[0.125in][r]{\thepage}%
    }%
  }%
}
%% end JL page numbers



%% Title
\title{Language agents achieve superhuman synthesis of scientific knowledge
%%%% Cite as
%%%% Update your official citation here when published 
% \thanks{\textit{\underline{Citation}}: 
% \textbf{Authors. Title. Pages.... DOI:000000/11111.}} 
}

\author{
\begin{tabular}{ccc}
Michael D. Skarlinski$^{1}$ & Sam Cox$^{1,2}$ & Jon M. Laurent$^{1}$ \\ \\
James D. Braza$^{1}$ & Michaela Hinks$^{1}$ & Michael J. Hammerling$^{1}$\\ \\
Manvitha Ponnapati$^{1}$ & Samuel G. Rodriques$^{1,3*}$ & Andrew D. White$^{1,2*}$
\end{tabular}
\\
\\
$^{1}$FutureHouse Inc., San Francisco, CA\\
$^{2}$University of Rochester, Rochester, NY\\
$^{3}$ Francis Crick Institute, London, UK\\
$^{*}$These authors jointly supervise technical work at FutureHouse.\\
Correspondence to: \texttt{\{sam,andrew\}@futurehouse.org}
}



\begin{document}
\maketitle

%\thispagestyle{firstpage}

\begin{abstract}
Language models are known to ``hallucinate'' incorrect information, and it is unclear if they are sufficiently accurate and reliable for use in scientific research. We developed a rigorous human-AI comparison methodology to evaluate language model agents on real-world literature search tasks covering information retrieval, summarization, and contradiction detection tasks. We show that PaperQA2, a frontier language model agent optimized for improved factuality, matches or exceeds subject matter expert performance on three realistic literature research tasks without any restrictions on humans (i.e., full access to internet, search tools, and time). PaperQA2 writes cited, Wikipedia-style summaries of scientific topics that are significantly more accurate than existing, human-written Wikipedia articles. We also introduce a hard benchmark for scientific literature research called LitQA2 that guided design of PaperQA2, leading to it exceeding human performance. Finally, we apply PaperQA2 to identify contradictions within the scientific literature, an important scientific task that is challenging for humans. PaperQA2 identifies $2.34 \pm 1.99$ (mean $\pm$ SD, $N=93$ papers) contradictions per paper in a random subset of biology papers, of which $70\%$ are validated by human experts. These results demonstrate that language model agents are now capable of exceeding domain experts across meaningful tasks on scientific literature. 
\end{abstract}
\vspace{0.7cm}

% keywords can be removed
% \keywords{AI \and Retrieval Augmented Generation \and Agents \and Benchmarks}
%\section{Teaser}

%\textbf{Short summary}: We develop novel human-AI comparison methods across scientific literature search tasks, and present an agentic system capable of superhuman performance on retrieval and summarization.

\begin{bibunit}[unsrt] 
\section{Introduction}
Large language models (LLMs) have the potential to assist scientists with retrieving, synthesizing, and summarizing the literature\cite{gao2023retrieval, shao2024assisting, Lo2023TheSR}, but still have several limitations for use in research tasks. Firstly, factuality is essential in scientific research, and LLMs hallucinate\cite{tonmoy2024comprehensive}, confidently stating information that is not grounded in any existing source or evidence. Secondly, science requires extreme attention to detail, and LLMs can overlook or misuse details when faced with challenging reasoning problems\cite{dahl2024large}. 
Finally, benchmarks for retrieval and reasoning across the scientific literature today are underdeveloped. They do not consider the entire literature, but instead are restricted to abstracts\cite{jin-etal-2019-pubmedqa}, retrieval on a fixed corpus\cite{krithara2023bioasq}, or simply provide the relevant paper directly\cite{jin2019pubmedqa}. These benchmarks are not suitable as performance proxies for real scientific research tasks, and, more importantly, often lack a direct comparison to human performance. Thus, it remains unclear whether language models and agents are suitable for use in scientific research.

%We believe these to be foundational tasks for scientists, and thus an LLM-based system that can exceed human performance on these tasks represents a significant milestone on the continued progress of LLMs in science. 

\begin{figure}
    \centering % Center the figure
    \includegraphics[width=0.8\textwidth]{Figures/Wikicrow-Schematic.png}
    \caption{\textbf{A}. Schematic of PaperQA2's agentic toolset along with relevant action representations within each tool. \textbf{B.} PaperQA2 performance across question answering, cited article summarization, and contradiction detection. Error bars represent standard error.} % Add a caption
    \label{fig:paperqa_summary} % Add a label for referencing the figure
\end{figure}

We therefore set out to develop a rigorous comparison between the performance of AI systems and humans on three real-world tasks: a retrieval task involving searching the entire literature to answer questions; a summarization task involving producing a cited, Wikipedia-style articles on scientific topics; and a contradiction-detection task, involving extracting all claims from papers and checking them for contradictions against all of literature. This is, to our knowledge, the first robust procedure for evaluating a single AI system on multiple real-world literature search tasks. Using our newly developed evaluations, we explored multiple designs leading to a system we call PaperQA2 (\textbf{\autoref{fig:paperqa_summary}A}), which exceeds the performance of PhD students and postdocs on the retrieval and summarization tasks. Applying PaperQA2 to the contradiction detection task enables us to identify contradictions in biology papers at scale (\textbf{\autoref{fig:paperqa_summary}B}). For example, a statement that the ZNF804A rs1344706 allele positively affects brain structure in schizophrenia patients \cite{wang2019pleiotropic} was found to be contradicted by a later publication which found that rs1344706's effects on cortical thickness, surface area, and cortical volume in the brain aggravate the risk of schizophrenia\cite{wei2015znf804a}.


\section{Answering scientific questions}

To evaluate AI systems on retrieval over the scientific literature, we first generated LitQA2,\cite{laurent2024lab} a set of 248 multiple choice questions with answers that require retrieval from scientific literature (\textbf{\autoref{fig:litqa_performance}A}). LitQA2 questions are designed to have answers that appear in the main body of a paper, but not in the abstract, and ideally appear only once in the set of all scientific literature. These constraints enable us to evaluate response accuracy by matching the system's cited source DOI with the DOI originally assigned by the question creator. To enforce these criteria, we generated large numbers of questions about obscure intermediate findings from very recent papers, and then excluded any questions where either an existing AI system or a human annotator could answer the question using an alternative source (\autoref{sec:question_evaluators}). These were generated entirely by experts, although there are emerging ideas about how to automate this process\cite{wan2024sciqag}. When answering LitQA2 questions, models can refuse to answer via selecting \textit{Insufficient information to answer this question}. Similar to prior work \cite{rajpurkar2018know} and matching actual scientific questions, some questions are intended to be unanswerable. We evaluate two metrics: precision, the fraction of questions answered correctly when a response is provided, and accuracy, the fraction of correct answers over all questions. We also consider recall, which is the total percentage of questions where the system attributed its answer to the correct source DOI denoted in LitQA2.

Having developed LitQA2, we then utilized it to design an AI system for the scientific literature. The current paradigm for eliciting factually-based responses from LLMs is to use retrieval-augmented generation (RAG)\cite{lewis2020retrieval, shuster-etal-2021-retrieval-augmentation}. RAG provides additional context to the LLM (e.g., snippets from research papers) to ground the generated response. As scientific literature is quite large, identifying the correct snippet is a challenge. Strategies like using metadata or hierarchical indexing can improve retrieval in this setting\cite{gao2023chatrec}, but finding the correct paper for a task often requires iterating and revising queries. Inspired by PaperQA\cite{lala2023paperqa}, PaperQA2 is a RAG \textit{agent} that treats retrieval and response generation as a multi-step agent task\cite{karpas2022mrkl} instead of a direct procedure. PaperQA2 decomposes RAG into \textit{tools}, allowing it to revise its search parameters and to generate and examine candidate answers before producing a final answer (\textbf{\autoref{fig:paperqa_summary}A}). PaperQA2 has access to a ``Paper Search'' tool, where the agent model transforms the user request into a keyword search that is used to identify candidate papers. The candidate papers are parsed into machine readable text, and chunked for later usage by the agent. PaperQA2 uses the state-of-the-art document parsing algorithm (Grobid\cite{GROBID}) that reliably parses sections, tables, and citations from papers. After finding candidates, PaperQA2 can use a ``Gather Evidence'' tool that first ranks paper chunks with a top-$k$ dense vector retrieval step, followed by an LLM reranking and contextual summarization (RCS) step. RCS prevents irrelevant chunks from appearing in the RAG context by summarizing and scoring the relevance of each chunk, which is known to be critical for RAG\cite{shi2023large}. The top ranked contextual summaries are stored in the agent's state for later steps. PaperQA2's design differs from similar RAG systems like Perplexity\cite{perplexity}, Elicit\cite{elicit}, or Mao et al.\cite{mao2020generation} which deliver retrieved chunks without substantial transformation in the context of the user query. While RCS is more costly than retrieval without a contextual summary, it allows PaperQA2 to examine much more text per user question. The RCS step also injects metadata about the source paper, like its citation count and journal. Once the PaperQA2 state has summaries, it can call a ``Generate Answer'' tool which uses the top ranked evidence summaries inside a prompt to an LLM for the final response to the asked questions or assigned task. To further improve recall, PaperQA2 adds a new ``Citation Traversal'' tool (\autoref{sec:cit-trav}) that exploits the citation graph as a form of hierarchical indexing to add additional relevant sources. 

%We were able to identify three major changes that improved the performance of PaperQA2 over PaperQA. Firstly, we developed a \textit{Citation Traversal} tool (\autoref{sec:cit-trav}) that exploits the citation graph; \textbf{MORE EDITING HERE: what was the impact of each change} (2) a state-of-the art document parsing algorithm (Grobid\cite{GROBID}) that reliably parses sections, tables, and citations; and (3) an improved RCS step via enhanced summary metadata and parameter optimization. %Utilizing these optimizations, we show that language agents like PaperQA2 are capable of achieving superhuman performance on retrieval and summarization tasks in a realistic scientific setting. We apply PaperQA to systematically detect contradictions across the scientific literature, and find across a random subset of biology papers \textbf{TODO HERE}

In answering LitQA2 questions, PaperQA2 parsed and utilized an average of 14.5 $\pm$ 0.6 (mean $\pm$ SD, $n=3$) papers per question. Running PaperQA2 on LitQA2 yielded a precision of 85.2\% $\pm$ 1.1\% (mean $\pm$ SD, $n=3$), and an accuracy of 66.0\% $\pm$ 1.2\% (mean $\pm$ SD, $n=3$) (\textbf{\autoref{fig:litqa_performance}B}), with the system choosing ``insufficient information'' in 21.9\% $\pm$ 0.9\% (mean $\pm$ SD, $n=3$) of answers. To compare the performance of PaperQA2 to other retrieval systems, we evaluated the performance of PaperQA with original parameters, commercial systems like Perplexity\cite{perplexity} and Elicit\cite{elicit}, and frontier (non-RAG models) on LitQA2. We found that PaperQA2 outperforms other RAG systems on the LitQA2 benchmark in both precision and accuracy. We also found that all RAG systems tested, with the exception of Elicit, outperform non-RAG frontier models in both precision and accuracy. 

To ensure that we did not overfit PaperQA2 to achieve high performance on LitQA2, we generated a new set of 101 LitQA2 questions after making most of the engineering changes to PaperQA2. The accuracy of PaperQA2 on the original set of 147 questions did not differ significantly from its accuracy on the latter set of 101 questions, indicating that our optimizations in the first stage generalized well to new and unseen LitQA2 questions (\autoref{tab:litqa_before_after}).

To compare PaperQA2 performance to human performance on the same task, human annotators who either possessed a PhD in biology or a related science, or who were enrolled in a PhD program (see \autoref{sec:question_evaluators}), were each provided a subset of LitQA2 questions and a performance-related financial incentive of \$3-12 per question to answer as many questions correctly as possible within approximately one week, using any online tools and paper access provided by their institutions. Under these conditions, human annotators achieved 73.8\% $\pm$ 9.6\% (mean $\pm$ SD, $n=9$) precision on LitQA2 and 67.7\% $\pm$ 11.9\% (mean $\pm$ SD, $n=9$) accuracy (\textbf{\autoref{fig:litqa_performance}A}, green line). PaperQA2 thus achieved superhuman precision on this task ($t(8.6)=3.49, p=0.0036$) and did not differ significantly from humans in accuracy ($t(8.5)=-0.42, p=0.66$).

\begin{figure}[h!]
    \centering
    \includegraphics[width=0.8\textwidth]{Figures/Wikicrow-PaperQA.png}
    \caption{\textbf{A.} Example LitQA2 question, PaperQA2 answer, and metadata. \textbf{B.} PaperQA2 performance on LitQA2 across technologies. \textbf{C.} PaperQA2 performance studies and ablations across component categories. Error bars are 95\% CI. \textbf{D.} Aggregated LitQA2 DOI recall per PaperQA2 stage. \textit{Search Recall} includes DOIs found via the ``Paper Search'' or ``Citation traversal'' tools, \textit{Top-k Ranking} includes all DOIs with similarity rankings below the top-k ranking cutoff (30), \textit{RCS} includes all DOIs selected by RCS, and \textit{Attribution} includes all DOIs cited in the ``Generate Answer'' tool.}
    \label{fig:litqa_performance} % Add a label for referencing the figure
\end{figure}

\section{Performance analysis of PaperQA2}
We varied the parameters of PaperQA2 to understand which are responsible for its accuracy (\textbf{\autoref{fig:litqa_performance}C}). We created a non-agentic version (\textit{No Agent}) which had a hard-coded sequence of actions (paper search, gather evidence, then generate answer). The non-agentic system had significantly lower accuracy ($t(3.7)=3.41, p=0.015$), validating the choice of using an agent. We attribute the performance difference to the agent's better recall because it can return to and change keyword searches (paper search tool calls) after observing the amount of relevant papers it finds. The highest accuracy LitQA2 runs had $1.26 \pm 0.07$ (mean $\pm$ SD) searches per question, and $0.46 \pm 0.02$ (mean $\pm$ SD) citation traversals per question showing that the agent will sometimes return to an additional search or traverse the citation graph to gather more papers.

We then explored the tools and their parameters available to PaperQA2. An important RAG parameter is the number of contexts or text chunks to include in the final text generation step (``Generate Answer'' tool in \textbf{\autoref{fig:paperqa_summary}A}). This is a balance because increasing the count improves the chance of including the key context needed to answer a question (improving accuracy), but also increases the amount of distracting irrelevant context that reduces precision\cite{liu2023lostmiddlelanguagemodels}. We varied the amount of contexts from 15 to 5 in \textbf{Figure S\ref{fig:litqa_precision}} and \textbf{\autoref{fig:paperqa_summary}} to see this effect; 15 gives highest precision and 5 gives highest accuracy. 

To improve relevant chunk retrieval, we hypothesized that papers found as either citers or citees of existing relevant chunks would be an effective form of hierarchical indexing. This was validated by ablating the ``Citation Traversal'' tool (\textit{No Cit. Trav.}), which showed an increased accuracy ($t(2.55)=2.14, p=0.069$), and significantly increased DOI recall ($t(3)=3.4$, $p=0.022$) at all stages of the PaperQA2 flow. (\textbf{\autoref{fig:paperqa_summary}D}) This tool's process mirrors the way that scientists interact with the literature.

To quantify retrieval accuracy changes across LLM implementations, we performed experiments which varied the model choice for our ``Generate Answer'' tool as well as the RCS step in our ``Gather Evidence'' tool. Both of these tools give LLMs the opportunity to correctly identify crucial information in our scientific corpus, and we wanted to evaluate if their combined usage would be more effective than an un-transformed insertion of the top chunks into the final context window. The \textit{No RCS Model} ablation validates that adding RCS to a traditional RAG text generation step significantly increases retrieval accuracy ($t(3.92)=9.29, p < 0.001$). Interestingly, this is not true across all models tested, smaller models (GPT-3.5-Turbo (RCS), Llama3 (70B)) decrease overall accuracy when used for RCS, relative to not using a model at all. This indicated there is a comprehension threshold that must be met for effective summarization and relevance evaluation. GPT-4-Turbo significantly outperformed other models on LitQA2 accuracy when used in the RCS step ($t(3.47)=6.14, p=0.003$). Finally, Claude-Opus \cite{claude2024} had the highest LitQA2 precision (\textbf{\autoref{fig:litqa_precision}}), though it was not a significant increase over Gemini-1.5-Pro and GPT-4-Turbo.

To reduce the large cost of generating contextual summaries, we examined the effect of lowering the top-k ranking depth, i.e. limiting the number of passages considered in the RCS step. We see a significant increase in accuracy with increasing depth (i.e. more document chunks entering into the RCS ranking) from 1 to 10 ($t(2.15)=5.44, p=0.014$), and a long tail of diminishing performance gains from 10 to 30 (the default). In summary, having a deep (>10) RCS ranking list, with a high-performing LLM, was crucial in achieving human-level accuracy on LitQA2. 

We had hypothesized that parsing quality would affect accuracy, but Grobid parsings and larger chunk sizes did not significantly increase precision, accuracy, or recall on LitQA2 (\autoref{fig:chunksize_ablations}). This is likely specific to being a retrieval task, as there is often only a single passage needed from a paper's body, which makes our result insensitive to parser changes. Anecdotally, we found better parsings to be crucial for extracting data from tables in WikiCrow (detailed in the \autoref{sec:si-wikicrow}).
%

\section{Summarizing scientific topics} \label{sec:wikicrow}

To evaluate PaperQA2 on summarization, we engineered a system called WikiCrow, which generates cited Wikipedia-style articles about human protein-coding genes by combining several PaperQA2 calls on topics such as the structure, function, interactions, and clinical significance of the gene (\textbf{\autoref{fig:WikiCrow_results}A}). There has been previous work on \textit{unconstrained} document summarization, where the document must be found and then summarized,\cite{Giorgi2022ExploringTC} and even writing Wikipedia-style articles with RAG\cite{shao2024assisting}. These studies have not compared directly against Wikipedia with human evaluation. Instead, they used either LLMs to judge or compared ROGUE (text overlap) against ground-truth summaries. Here, we measure directly against human-generated Wikipedia with subject mater expert grading.

We used WikiCrow to generate 240 articles on genes that already have non-stub Wikipedia articles to have matched comparisons. WikiCrow articles averaged 1219.0 $\pm$ 275.0 words (mean $\pm$ SD, $N=240$), longer than the corresponding Wikipedia articles (889.6 $\pm$ 715.3 words). The average article was generated in 491.5 $\pm$ 324.0 seconds, and had an average cost of \$4.48 $\pm$ \$1.02 per article (including costs for search and LLM APIs). We compared WikiCrow and Wikipedia on 375 statements sampled from the 240 paired articles. Statements were selected using cues from document formatting (\autoref{sec:si-wikicrow}). The initial article sampling excluded any Wikipedia articles that were ``stubs'' or incomplete articles. Statements were then shuffled and given, blinded, to human experts, who graded statements according to whether they were (1) cited and supported; (2) missing a citation; or (3) cited and unsupported. We found that WikiCrow had significantly fewer ``cited and unsupported'' statements than the paired Wikipedia articles (13.5\% vs. 24.9\%) ($p=0.0075, \chi^2(1), N=375$ for all tests in this section). WikiCrow failed to cite sources at a 3.9x lower rate than human written articles, as only 3.5\% of WikiCrow statements were uncited, vs. 13.6\% for Wikipedia ($p < 0.001$). In addition, defining precision for WikiCrow as the ratio of cited and supported statements over all cited statements, we found that WikiCrow displayed significantly higher precision than human-written articles (86.1\% vs. 71.2\%, $p = 0.0013$).

\begin{figure}[h!]
    \centering % Center the figure
    \includegraphics[width=0.8\textwidth]{Figures/Wikicrow-WikiArticle.png}
    \caption{\textbf{A.} WikiCrow prompt graph, showing the initial seed variable (gene name), PaperQA2 prompts for each section, and overview LLM call. Each section feeds into a Python script which stitches the outputs together. \textbf{B.} Sample WikiCrow article for gene FAM83H, with truncated sections under each header. \textbf{C.} Performance statistics for a WikiCrow vs. Wikipedia comparison performed by evaluators. \textbf{D.} ``Cited and unsupported'' issue categorization counts.}
    \label{fig:WikiCrow_results} % Add a label for referencing the figure
\end{figure}

The ``cited and unsupported'' evaluation category includes both inaccurate statements (e.g. true hallucinations or reasoning errors) and statements that are accurate with inappropriate citations. To investigate the nature of the errors in Wikipedia and WikiCrow further, we manually inspected all reported errors and attempted to classify the issues as follows: reasoning issues, i.e. the written information contradicts, over-extrapolates, or is unsupported by any included citations; attribution issues, i.e. the information is likely supported by another included source, but either the statement does not include the correct citation locally or the source is too broad (e.g. a database portal link); or trivial statements, which are true passages, but overly pedantic or unnecessary (\textbf{\autoref{fig:WikiCrow_results}D}). Surprisingly, we found that compared to Wikipedia, WikiCrow had significantly fewer reasoning errors (12 vs. 26, $p = 0.0144, \chi^2(1), N=375$) but a similar number of attribution errors (10 vs. 16, $p = 0.21$), suggesting that the improved factuality of WikiCrow over Wikipedia was largely due to improvements in reasoning. Although language models are clearly prone to reasoning errors (or hallucinations), in our task at least they appear to be less prone to such errors than Wikipedia authors or editors. This statement is specific to the agentic RAG setting presented here: language models like GPT-4 on their own, if asked to generate Wikipedia articles, would still be expected to hallucinate at high rates.
\footnote{All WikiCrow generated articles for this study are now available for download from a Google Cloud Storage bucket \texttt{https://storage.googleapis.com/fh-public/wikicrow2/}. A command line tool like \texttt{gsutil} can be used to list and bulk-access these files. All WikiCrow and LitQA evaluator responses are available via the following \href{https://docs.google.com/spreadsheets/d/1QvfxHSplVb-YyyDBkko_mhmbcomfOjBNtuW06f87jd0/edit?usp=sharing}{link}.}


\section{Detecting contradictions in the literature}

Because PaperQA2 can explore scientific literature at much higher throughput than human scientists, we reasoned that we could deploy it to systematically identify contradictions and inconsistencies in the literature at scale. Contradiction detection is a ``one versus many'' problem, which in principle involves comparing claims or statements in one paper with all other claims or statements in the literature. At scale, contradiction detection becomes a ``many versus many'' problem and loses feasibility for humans. Thus, we leveraged PaperQA2 to build a system called ContraCrow that automatically detects contradictions in the literature (\textbf{\autoref{fig:contradiction_detection}A}). 

Contradiction detection is also known as claim verification or colloquially as ``fact checking''\cite{vlachos2014fact}. This task has been studied for over a decade, especially in the context of claims in the news or the internet\cite{vlachos2014fact, ferreira2016emergent}. Although originally restricted to context and a claim, the setting extended to be unconstrained\cite{popat2017truth} and recent work tries to work at the scale of the internet\cite{schlichtkrull2024averitec}. Some claim verification work has also focused on scientific claims\cite{wadden2020fact, wadden2021multivers}. The main novelty of this work is in detecting contradictions without a restricted corpus and evaluating with human experts, not against a benchmark. 

ContraCrow first extracts claims from a provided paper using a series of LLM completion calls (similar to Schlichtkrull et al. (2024)\cite{schlichtkrull2024averitec}), and then feeds those claims into PaperQA2 with a \textit{contradiction detection} prompt. This prompt instructs the system to evaluate whether there are any contradictions in literature to the provided claim, providing both an answer and a choice from an 11-point Likert scale (\textbf{\autoref{fig:contradiction_detection}B}, Methods \autoref{sec:cc-methods}). Utilizing a Likert Scale allows the system to give more reliable and interpretable scores when providing rank\cite{abeysinghe2024challenges}.

\begin{figure}[h!]
    \centering
    \includegraphics[width=0.8\textwidth]{Figures/Wikicrow-Contradictions.png}
    \caption{\textbf{A.} Schematic of ContraCrow. \textbf{B.} Likert scale used to evaluate contradictions, along with the integer mapping. \textbf{C.} ROC curve and metrics (Likert scale threshold of 8) of ContraCrow's performance on ContraDetect benchmark. \textbf{D.} Bar plot showing percentage of all claims found in 93 papers that are given scores of 8, 9, and 10 by ContraCrow. \textbf{E.} Histogram of the number of contradictions found per paper over 93 papers. The table shows the average number of claims ($35.16 \pm 21.72$ [mean $\pm$ SD, $N=93$]) and contradictions ($2.34 \pm 1.99$ [mean $\pm$ SD, $N=93$]) per paper found. \textbf{F.} Proportion of ContraCrow contradictions with scores 8 and 9-10 validated by expert evaluators. \textbf{G.} Range of evaluator scores for each claim. Bars represent the range of Likert scale scores assigned by expert evaluators and stars represent the ContraCrow label on the same claims. The table shows a mean standard deviation per claim of 1.45 ($N=30$ claims) over all evaluators and a mean Pearson correlation coefficient of $0.579 \pm 0.197$ (mean $\pm$ SD) between evaluators ($N=5$).}
    \label{fig:contradiction_detection}
\end{figure}

To evaluate ContraCrow, we first derived a contradiction detection benchmark, ContraDetect, from LitQA2, as detailed in \autoref{sec:cc-eval}. Briefly, we converted half of the question-answer pairs in LitQA2 into declarative, incorrect statements that are contradicted by the papers referenced in the corresponding LitQA2 question. (For example, a question like ``what color is grass?'' would become ``grass is purple.'') We converted the other half of the LitQA2 questions into declarative, correct statements that are supported by the corresponding papers (``what color is grass'' becomes ``grass is green.''). 

We then evaluated ContraCrow on its ability to detect the contradictions in ContraDetect. By transforming the Likert scale output into integers, we were able to tune the detection threshold and obtain an ROC curve with an AUC of $0.842$ (\textbf{\autoref{fig:contradiction_detection}C}). Setting a threshold of 8 (\textit{contradiction}), ContraCrow achieved $73\%$ accuracy, $88\%$ precision, and a false positive rate of only $7\%$. To further evaluate precision and to demonstrate ContraCrow's ability to handle ``Has anybody ever done x?'' questions, we then ran ContraCrow on 42 ``no-evidence statements'' that have never to our knowledge been reported on (detailed in \autoref{sec:cc-eval}). These claims were hand-generated by the authors within their fields of expertise, and ContraCrow correctly chose \textit{lack of evidence} (5) as its response $98\%$ of the time, indicating its ability to distinguish between real contradictions and lack of support.

We then applied ContraCrow to a set of 93 biology-related papers randomly selected from our database, identifying an average of $35.16 \pm 21.72$ (mean $\pm$ SD, $N=93$) claims per paper. Of the 3,180 claims analyzed over the 93 papers, $6.85\%$ were deemed by ContraCrow to be contradicted by the literature, with $2.89\%$, $3.77\%,$ and $0.19\%$ assigned scores $8$, $9$, and $10$, respectively (\textbf{Figure \ref{fig:contradiction_detection}D}). Setting a Likert scale threshold of $8$, we detected an average of $2.34 \pm 1.99$ contradictions per paper (mean $\pm$ SD) (\textbf{Figure \ref{fig:contradiction_detection}E}). As an example, one contradiction detected by our system concerned the prognostic implications of LEF1 expression in colorectal carcinomas. The source paper, Kriegl et al. (2010)\cite{kriegl2010lef}, finds using immunohistochemistry on a tissue microarray that LEF-1 correlates positively with longer overall survival. By contrast, a study published the following year also using tissue microarrays found LEF1 overexpression in colorectal cancer correlates negatively with longer overall survival and is also correlated with liver metastasis\cite{lin2011comparative}, results that have been supported by other studies as well\cite{wang2013increased,wang2013knockdown,kermanshahi2014lef}, thus explicitly contradicting the original statement.

To evaluate the validity of the contradictions detected this way, expert human annotators evaluated 50 claims that were assigned ContraCrow scores of 8 and 50 claims that were assigned scores of 9 or 10, considering all evidence and reasoning cited by the model. Annotators agreed with ContraCrow's findings on $70\%$ of evaluated claims, or 1.64 contradictions per paper, demonstrating significant agreement ($p = 1e^{-4}, \chi^2(1)$), with an F1 score of $0.82$. Interestingly, claims scored as 8 by ContraCrow were not more likely to elicit human agreement (70\%) than claims scored as 9 or 10 (70\%) (\textbf{Figure \ref{fig:contradiction_detection}F}). The final number, $1.64$ validated contradictions per paper, serves as a lower bound on the abundance of human-validatable contradictions in the biology literature.

We were concerned that the annotators on the ``contradiction validation'' task might exhibit some bias toward agreeing with the model or otherwise be influenced by the model's reasoning and chosen sources. We therefore further evaluated ContraCrow by asking human experts to identify contradictions in the same literature considered by ContraCrow, without any access to ContraCrow's reasoning. This ``contradiction detection'' task (\autoref{sec:cc-eval}) is a much more challenging task for annotators than the contradiction validation task, because it requires them to exhaustively consider a significant amount of literature relevant to the claim in question, whereas the contradiction validation task only requires them to consider the specific sections cited by the model. In the contradiction detection task, human experts were provided with the top 15 paper chunks identified as relevant to the claim by PaperQA2 (the same chunks ContraCrow had access to), and were asked to respond whether the claim was contradicted by the available evidence on both the Likert scale and binary determination (``yes'' or ``no''). Experts considered a mixture of claims rated as contradictions and non-contradictions by ContraCrow. We found that experts' binary responses agree with each other on $75.5\% \pm 13.43\%$ of claims provided (mean $\pm$ SD, $N=10$ pairwise comparisons between human annotators), whereas they agree with ContraCrow on $60.42\% \pm 5.99 \%$ of claims (mean $\pm$ SD, $N=5$ human annotators), indicating that humans are significantly more correlated with each other than they are with ContraCrow ($p=0.015$). Examining the Likert scale values over 30 claims evaluated by human annotators leads us to speculate that overconfidence on ContraCrow's part is the primary driver of its lack of agreement with human annotators, and indicates directions for future improvement (\textbf{Figure \ref{fig:contradiction_detection}G}).

We are thus able to establish a lower bound for the average number of human-validatable contradictions per paper in biology literature at 1.64. Importantly, just because a claim is detected as a contradiction does not mean that the claim is wrong. For example, one detected contradiction is \textit{``GBP is only found in the cytosol of human fibroblasts.''}\cite{cheng1985affinity} This example shows how scientific literature and findings can update and contradict over time, as ContraCrow points out: \textit{``...more recent research has shown that GBPs can localize to various cellular compartments and membranes\cite{britzen2010intracellular, modiano2005golgi, naschberger2006human, schelle2023functional}.''} Thus, claims can contradict but still be contextually valid, and contradictory claims may be merely reflective of the iterative nature of research.

\section{Conclusions}

We developed a methodology to compare or validate AI systems against human performance in realistic tasks for scientific research. PaperQA2 outperforms human experts on answering questions across all scientific literature; produces summaries that are, on average, more factual than Wikipedia summaries; and can be deployed to identify contradictions in scientific literature at scale. The contradiction work, in particular, attests to the potential of systems such as PaperQA2 for science: notably, one human expert who performed this task, who was also tested on a large battery of other benchmarks, reported without solicitation that the contradiction detection task was the hardest task they were asked to perform. Although PaperQA2 is expensive compared to lower accuracy commercial systems, it is inexpensive in absolute terms, costing \$1 to \$3 per query. Scaling up PaperQA2 and other literature-enabled agents like WikiCrow and ContraCrow empowers us to take advantage of the latent knowledge in literature at much greater scale than is possible today.


\section{Data Availability}

The code necessary to replicate these results or modify the algorithm for further research is included on Github via \href{https://github.com/Future-House/paper-qa}{paperqa}. Data including all evaluator responses, contradiction detection claims, litQA questions, and WikiCrow candidate statements are available in the supplementary materials. All generated WikiCrow articles for this study are available in a public Google Cloud bucket here: \texttt{https://storage.googleapis.com/fh-public/wikicrow2/}.

\paragraph{Author Contributions} MDS led work on PaperQA2 and WikiCrow, in collaboration with JDB. SC led work on ContraCrow. JML led work on LitQA. MH, MJH, and MP contributed questions to the LitQA benchmark. SGR and ADW conceived of and supervised the project.

\paragraph{Acknowledgments} Work at FutureHouse is supported by the generosity of Eric and Wendy Schmidt. We also acknowledge Matt Rubashkin for contributions to an earlier stage of the project, and we acknowledge all current members of FutureHouse for useful discussions, including Sid Narayanan, Ryan Rhys-Griffiths, Cade Gordon, Peter Chang, and Conor Igoe. We also acknowledge support and resources from the Semantic Scholar Project at the Allen Institute for AI.

\paragraph{Competing interests} The authors report no competing interests. FutureHouse, Inc. is a non-profit research organization. 

\putbib[references] % Include the main

\end{bibunit}

\newcounter{suppfigure}
\renewcommand{\thesuppfigure}{S\arabic{suppfigure}}
\setcounter{suppfigure}{0}

\section{Methods}

\subsection{PaperQA Implementation and Parameters}

All reported figures and data in this work were built on the open source PaperQA package, available on GitHub at \href{https://github.com/Future-House/paper-qa}{paperqa}. While the core PaperQA repository provides the basic algorithms used, it does not include the Grobid parsing code, access to non-local full-text literature searches, or the citation traversal tool. The open source version of PaperQA utilizes LangChain \cite{LANGCHAIN} for its agentic and state update operations. The full configuration objects for all experiments run in this paper are included for further customization.

Note that while \texttt{paperqa} gives the ability to recreate this work, the experiments reported in this paper were performed using a more featureful HTTP server that takes advantage of bespoke infrastructure at the authors' institution. This infrastructure includes features such as user authentication, MongoDB request caching, Redis object caching and global-load balancing, several PostgreSQL DBs with associated ORM code, cost-monitoring modules, time-profiling modules, configuration storage and run orchestration (Dagster\cite{dagster} and kubernetes\cite{kubernetes}), cloud bucket storage for PDFs, a CI pipeline with semi-automated deployments, and infrastructure code for deploying auto-scaling instances in the cloud. None of these features affect performance on a per-query basis, but provide increased scalability, measurability, and persistence. To run our same server infrastructure, users would need to provision all of these assets and configure the deployments themselves. \texttt{paperqa} should serve allow usage and customization sufficient for most research purposes, and should be sufficient to reproduce the results reported here.

Even within \texttt{paperqa}, the ``Paper Search'' tool is limited by access to full text repositories of scientific papers, often bound by licensing agreements. The included implementation only works from local files accessible to each user. Our implementation starts with a full or partial-text keyword search, where the keywords have been specified by the agent when selecting the paper search tool. The ranked results returned from these services are then matched to a user's existing paper repository or can be retrieved on-the-fly if open-access or partner links exist for these works. These matching papers are parsed, and pulled into our agent state for usage with other tools. Note that the search services will have access to a larger corpus of works than is available to us via our repository and accessible link traversal, in these cases the system will simply skip these papers and they are not used. A stub of the paper search tool is implemented in \texttt{paperqa} with directions for users to implement their own retrieval since it will be limited to their own access to full-text papers. 

Ablations and configurations for workflows like WikiCrow are exposed in \texttt{paperqa} as nested configuration objects. All experiments performed in this work correspond to included configuration objects. Here we highlight the configuration variable descriptions corresponding to the salient features tested in this work, though all variable names are available in the included files. 

\begin{itemize}
\item \texttt{query}: The main query task asked of the PaperQA agent, i.e. a LitQA question or a directive to write an article. 
\item \texttt{llm}: The LLM used in the generate answer tool can be a valid Anthropic, OpenAI, or Gemini model identifier. This parameter was varied for the model experiments in \autoref{fig:litqa_performance}.
\item \texttt{agent\_llm}: The LLM used for the agent orchestration, in this work, it was always fixed to \texttt{gpt-4-turbo-2024-04-09}.
\item \texttt{summary\_llm}: The LLM used for the RCS step in the gather evidence tool, must be a valid Anthropic, OpenAI, or Gemini model identifier. This parameter was varied for the model ablations in \autoref{fig:litqa_performance}.
\item \texttt{prompts}: A \texttt{PromptCollection} object from PaperQA \cite{lala2023paperqa}, which allows for specification of prompts in each tool, as well as features like turning off RCS (via \texttt{prompts.skip\_summarization}). The \textit{No RCS Model} ablation used this input as well as the WikiCrow prompts.
\item \texttt{max\_sources}: The number of top ranked sources to be included in the generate answer tool, in \textbf{\autoref{fig:paperqa_summary}A}, the `filter top summaries'' cutoff. This parameter was 5 in the top-performing \textit{Answer cutoff @ 5}, but 15 for all other experiments.
\item \texttt{consider\_sources}: The top-k cutoff, i.e. the number of chunks that will be used in the RCS step. This parameter was set to 30 by default in LitQA experiments, save for the \textit{Top-k Rank @ X} experiment where it was set to X. Additionally, for our WikiCrow prompts this parameter was set to 25. 
\item \texttt{agent\_tools}: An ordered list of tool names that will be used by the agent, including \texttt{gather\_evidence}, \texttt{paper\_search}, \texttt{generate\_answer}, and \texttt{citations\_traversal}. This always included all four tools except for the \textit{No Cit. Trav.} and \textit{No Agent} runs where \texttt{citations\_traversal} was excluded.
\item \texttt{docs\_index\_mmr\_lambda}: A pre-gather evidence MMR lambda parameter which can be used to pre-filter similar papers by name before gathering evidence. This was set to 0.9 for our WikiCrow run to promote diversity of sources, but 1.0 for LitQA experiments. 
\item \texttt{parsing\_configuration.ordered\_parser\_preferences}: A list of the parsing algorithm to use, either \texttt{paperqa\_default} (PyMuPDF) or \texttt{grobid}. \texttt{paperqa\_default} was the default for each ablation, and \texttt{grobid} was used for WikiCrow generation. This parameter was also varied in the experiments shown in \textbf{Figure \ref{fig:chunksize_ablations}}. 
\item \texttt{parsing\_configuration.chunksize}: the chunk size (in characters) to be used when chunking parsed documents. This parameter was varied in the experiments shown in \textbf{Figure \ref{fig:chunksize_ablations}}.
\item \texttt{parsing\_configuration.overlap}: the overlap (in characters) that will be common between sequential chunks. This was fixed at 750 for this work. 
\item \texttt{parsing\_configuration.chunking\_algorithm}: the algorithm used to chunk documents, \texttt{simple\_overlap} simply uses a sliding window with overlap, and \texttt{sections} uses semantic parsing by section (i.e. one chunk per section where possible), if sections need to be broken into multiple chunks the system will automatically handle this. \texttt{sections} is only supported via \texttt{parsing\_configuration.ordered\_parser\_preferences}=\texttt{grobid}. This parameter was varied in the experiments shown in {Figure \ref{fig:chunksize_ablations}}, and in our WikiCrow generation.
\item \texttt{temperature}: temperature used for the LLM in the generate answer tool.  This was set to 0 for all runs in this work.
\item \texttt{summary\_temperature}: temperature used for the LLM in the gather evidence tool's RCS step, this was set to 0 for all runs in this work.
\end{itemize}

\subsubsection{Tool implementations} \label{sec:cit-trav}

PaperQA2's agentic tools were implemented as in PaperQA\cite{lala2023paperqa}. Our agent was prompted with the following message to guide tool usage: 

\begin{quote}
Answer question: \{question\}. Search for papers, gather evidence, collect papers cited in evidence then re-gather evidence, and answer. Gathering evidence will do nothing if you have not done a new search or collected new papers. If you do not have enough evidence to generate a good answer, you can:
\\
- Search for more papers (preferred)\\
- Collect papers cited by previous evidence (preferred)\\
- Gather more evidence using a different phrase\\

If you search for more papers or collect new papers cited by previous evidence, remember to gather evidence again. Once you have five or more pieces of evidence from multiple sources, or you have tried a few times, call \{gen\_answer\_tool\_name\} tool. The \{gen\_answer\_tool\_name\} tool output is visible to the user, so you do not need to restate the answer and can simply terminate if the answer looks sufficient. The current status of evidence/papers/cost is \{status\}
\end{quote}

Where variables like \{status\} are included to represent the current state to the agent. Tools were implemented with the following prompts and settings. \\

\textbf{Paper Search Tool} \\

The paper search tool uses an initial keyword search, generated by the agent in the context of the user query. The agent is prompted as follows: 

\begin{quote}
A search query in this format: [query], [start year]-[end year]. You may include years as the last word in the query, e.g. 'machine learning 2020' or 'machine learning 2010-2020'. The current year is \{get\_year()\}. The query portion can be a specific phrase, complete sentence, or general keywords, e.g. 'machine learning for immunology'.
\end{quote}

Our initial search relies on services like Semantic Scholar\cite{semantic-scholar}, where candidate lists (default of 12) of relevant papers are generated then parsed. When parsed, the papers are first turned into text using either Grobid or PyMuPDF, then split into \texttt{chunksize} character sized pieces. If the \texttt{sections} parsing is used, then section chunks are split on header metadata provided by Grobid. An embedding vector is generated for each chunk using a hybrid implementation which concatenates a dense and sparse, keyword based embedding model. For the experiments included in this study, OpenAI's \texttt{text-embedding-large-3} was used. It was concatenated with a normalized 256 dimension vector which used modulus-encoding to extract a hot-encoded keyword from the tokenization integers provided by OpenAI's tiktoken \cite{tiktoken} library. These text chunks are put into a document context which is accessible by the agent for further manipulation with tools. The PaperQA entrypoint for these functions can be found on \href{https://github.com/Future-House/paper-qa}{github}. \\ \\


\textbf{Gather Evidence Tool} \\

As detailed in \href{https://github.com/Future-House/paper-qa}{github} for PaperQA, the Gather Evidence tool begins with a top-k vector ranking step, using the embedding vectors created in the Paper Search tool. The user query is embedded with the same algorithm, and cosine similarity is used to match all document chunks in the agent's context with the user query. The top-k chunks are then selected for the RCS step. 

The reranking and contextual summarization step most differentiates PaperQA's implementation relative to other RAG technologies. The tool's prior step, an top-k vector retrieval ranking, is a widely implemented \cite{ma2023queryrewritingretrievalaugmentedlarge, gao2023retrieval} approach to identify relevant documents, however, the RCS second step, is unique to PaperQA (to the authors' knowledge). While performance improvements with both deep reranking (or LLM) models and map-reduced summarizations \cite{sarthi2024raptor, zhuang2023opensourcellm, gao2023chatrec, yu2024rankrag} are well documented, combining the reranking operation with a contextual summary provides novel benefits. 

The step is implemented by mapping an LLM completion across each top-k chunk (system prompt):
\begin{quote}
Provide a summary of the relevant information that could help answer the question based on the excerpt. The excerpt may be irrelevant.  Do not directly answer the question - only summarize relevant information. Respond with the following JSON format:
\{\{ "summary": "...",  \\
"relevance\_score": "..."
\}\}
where "summary" is relevant information from text - \{summary\_length\} words and "relevance\_score" is the relevance of "summary" to answer the question (integer out of 10)
\end{quote}

Where each chunk is injected as follows:
\begin{quote}
    Excerpt from {citation}
    ----
    \{text\}
    ----
    Query: \{question\}
\end{quote}

After completion, each JSON object is parsed and the passages are re-ranked according to the new relevance scores. When running with WikiCrow, gene names are also prompted to be extracted as additional JOSN keys, these are kept and injected in the final answering context. Advantages of the RCS step are as follows: 1. Token usage efficiency is vastly improved, a contextual summary will be 200-400 tokens compared with our standard document's chunk size of 2,250 tokens. This allows for a significantly more accessible document corpus for injection into PaperQA's answering context window. Furthermore, we see no decrease in summarization efficacy, using LitQA performance as a proxy, across document chunk sizes from 750-3,000 tokens. 2. As a new feature in this work, the LLM can be prompted to provide its summary in a structured JSON or XML format to simplify its downstream data extraction. In addition to a relevance score used for reranking, this structure can include metadata (such as a gene name) which will be retained through the PaperQA workflow. This is used to reduce hallucination and confusion in the final answer context. Since the RCS step is performed in an embarrassingly parallel fashion, it's highly efficient, and its utility can be applied to an arbitrarily deep ranking, up to the rate or cost limits of the LLM API. Our studies on the efficacy of the RCS depth led us to use a much deeper RCS depth, and to utilize the best performing model for the RCS operation. This differs from the intuition in prior work \cite{lala2023paperqa}, which utilized a cheaper model during the RCS step. \\

\textbf{Generate Answer Tool} \\

This tool answers questions by taking a subset of the top ranked sources (from the RCS ranking), and injects them into a final context for answering. The default in this study was to inject 15 contextual summaries, but we saw maximal accuracy with 5 at the cost of precision. LLMs were prompted to answer as follows:

\begin{quote}
Answer the question below with the context.

Context:
\{context\}
----
Question: \{question\}

Write an answer based on the context. If the context provides insufficient information and the question cannot be directly answered, reply "I cannot answer." For each part of your answer, indicate which sources most support it via citation keys at the end of sentences, like (Example2012Example pages 3-4). Only cite from the context and only use the valid keys. Write in the style of a Wikipedia article, with concise sentences and coherent paragraphs. The context comes from a variety of sources and is only a summary, so there may inaccuracies or ambiguities. If quotes are present and relevant, use them in the answer. This answer will go directly onto Wikipedia, so do not add any extraneous information.

Answer (\{answer\_length\}):
\end{quote}

Where contexts are injected by the generate answer code before output is returned to the agent. \\

\textbf{Citation Traversal Tool} \\

Atop the PaperQA\cite{lala2023paperqa} tools, we created an additional tool to traverse one degree of citations, both forward in time (``future citers'') and backwards in time (``past references''). This tool enables a fine-grained search around paper(s) containing relevant information. The traversal originates from any paper containing a highly-scored contextual summary (RCS score 0-10), and our minimum score threshold was eight (inclusive). The papers corresponding to highly-scored summaries are referred to as $D_\text{prev}$ in \textbf{Algorithm \ref{alg:cit-trav}}. See \textbf{\autoref{tab:cit-trav-starting-dist}} for the frequencies of various $|D_\text{prev}|$ when this tool was selected.

\begin{table}[htb]
\caption{Various statistics on citation traversal.}
\begin{subtable}{\textwidth}
\centering
\begin{tabular}{| c || c c c c c c c c c c |}
\hline
$|D_\text{prev}|$ & 1 & 2 & 3 & 4 & 5 & 6 & 7 & 8 & 9 & 10 \\
\hline
Count & 2147 & 941 & 530 & 386 & 307 & 216 & 154 & 67 & 29 & 12 \\
\hline
Frequency (\%) & 44.8 & 19.6 & 11.1 & 8.1 & 6.4 & 4.5 & 3.2 & 1.4 & 0.6 & 0.3 \\
\hline
\end{tabular}
\caption{Distribution of traversal starting paper count $|D_\text{prev}|$.}
\label{tab:cit-trav-starting-dist}
\end{subtable}

\begin{subtable}{\textwidth}
\centering
\begin{tabular}{| c || c c c c c c c |}
\hline
$|D_\text{prev}|$ & 1 & 2 & 3 & 4 & 5 & 6 & 7+ \\
\hline
Frequency of 1 Overlap (\%) & \textbf{100.0} & \textbf{91.3} & \textbf{91.7} & 90.8 & 90.8 & 88.7 & 87.5 \\
\hline
Frequency of 2 Overlaps (\%) & \cellcolor{gray!25} & \textbf{8.7} & \textbf{7.4} & \textbf{7.5} & \textbf{7.2} & \textbf{8.2} & 8.5 \\
\hline
Frequency of 3 Overlaps (\%) & \cellcolor{gray!25} & \cellcolor{gray!25} & \textbf{0.9} & \textbf{1.4} & \textbf{1.4} & \textbf{2.0} & \textbf{2.4} \\
\hline
Frequency of 4 Overlaps (\%) & \cellcolor{gray!25} & \cellcolor{gray!25} & \cellcolor{gray!25} & \textbf{0.3} & \textbf{0.5} & \textbf{0.7} & \textbf{0.9} \\
\hline
Frequency of 5 Overlaps (\%) & \cellcolor{gray!25} & \cellcolor{gray!25} & \cellcolor{gray!25} & \cellcolor{gray!25} & \textbf{0.1} & \textbf{0.3} & \textbf{0.4} \\
\hline
Frequency of 6 Overlaps (\%) & \cellcolor{gray!25} & \cellcolor{gray!25} & \cellcolor{gray!25} & \cellcolor{gray!25} & \cellcolor{gray!25} & \textbf{0.1} & \textbf{0.2} \\
\hline
Frequency of 7+ Overlaps (\%) & \cellcolor{gray!25} & \cellcolor{gray!25} & \cellcolor{gray!25} & \cellcolor{gray!25} & \cellcolor{gray!25} & \cellcolor{gray!25} & \textbf{0.1} \\
\hline
\end{tabular}
\caption{Table showing the frequencies of citation overlap $o$ seen in LitQA, illustrating the percentage of traversed citations at stake when filtering with an overlap threshold $\theta_o$ . We chose to specify $\theta_o = \lceil \alpha \times |D_\text{prev}| \rceil$, where $\alpha$ is known as the overlap fraction and was defaulted to $\frac{1}{3}$. The bolded values show what overlaps would have been preserved using an $\alpha = \frac{1}{3}$.}
\label{tab:cit-trav-overlap-dist}
\end{subtable}
\end{table}

To first acquire citations, Semantic Scholar \cite{semantic-scholar} and Crossref \cite{crossref} APIs are called for past references and Semantic Scholar APIs are called for future citers. To collect all citations for a given paper, we make one API call per provider per direction, totalling four API calls/paper. All three providers only provide partial paper details, meaning a large fraction of the time a title or DOI is not present in the response metadata. To merge citations across providers, a best-effort de-duplication is performed using casefolded title and lowercased DOI. In \textbf{Algorithm \ref{alg:cit-trav}}, this logic takes place inside the \texttt{GetCitations} procedure.

Once citations have been acquired, bins of overlap $\mathcal{B}$ are computed. For example, traversing past references for the following six DOIs:
% Note: added a space in third one to enable line wrap
\texttt{10.1016/j.mcn.2006.08.007}, \texttt{10.1002/cpsc.17}, \texttt{10.1002/(sici)1098- 1136(200004)30:2<105::aid-glia1>3.0.co;2-h}, \texttt{10.1089/scd.2015.0244}, \texttt{10.1002/glia.22882}, and \texttt{10.1042/an20120041},
leads to one DOI cited by four papers, five DOIs cited by three papers, 29 DOIs cited by two papers, and 428 DOIs cited by just one paper.

To filter bins of overlap, a hyperparameter ``overlap fraction'' $\alpha$ was introduced to compute a threshold overlap $\theta_o$ as a function of the number of source papers ($|D_\text{prev}|$). For example, with an $\alpha = \frac{2}{5}$ and traversing from six source DOIs, all citations not appearing in at least three source DOIs were discarded. The default overlap fraction used in data collection was $\frac{1}{3}$. See \textbf{\autoref{tab:cit-trav-overlap-dist}} for a full distribution of overlaps seen during LitQA runs. Furthermore, a twelve paper limit $\ell$ was posed on the traversal, which meant in the above example only keeping six of the bin of 29 DOIs cited by two papers. To filter within a bin, we fall back on the count of future citers. This winnowing logic is detailed across \textbf{Algorithm \ref{alg:cit-trav}}'s \texttt{FilterOverlap} and \texttt{TraverseCitations} procedures. Lastly, we traverse both future citers and past references, feathering together the resultant DOIs before finding them. 

% An future optimization is traversing only citations from the source chunk of a contextual summary, as opposed to the entire paper. A comprehensive hyperparameter search over traversal direction (not both future/past), overlap fraction, and contextual summary score threshold could also be performed.

\begin{algorithm}
\setstretch{1.15}
\caption{Traverse Citations}
\label{alg:cit-trav}
\begin{algorithmic}[1]

\Require{Set of summaries $S$, score threshold $\theta_\text{score}$, overlap fraction $\alpha$, look future flag $\mathds{1}_\text{fut}$, limit $\ell$}
\Ensure{Set of traversed papers $D_\text{out}$, where papers are future citers if $\mathds{1}_\text{fut}$ else past references}
\Procedure{TraverseCitations}{$S$, $\theta_\text{score}$, $\alpha$, $\mathds{1}_\text{fut}$, $\ell$}
\State $D_\text{prev} \leftarrow \{s_d \mid s \in S \wedge s_\text{score} \geq \theta_\text{score} \}$
\Comment{Traverse from highly-scored summaries' corresponding papers}
\State $\mathcal{D} \leftarrow \Call{GetCitations}{D_\text{prev}, \mathds{1}_\text{fut}}$
\Comment{$\mathcal{D}$ is a set of sets of papers such that $|\mathcal{D}| = |D_\text{prev}|$}
\State $\theta_o \leftarrow \lceil \alpha \times |\mathcal{D} | \rceil$
\Comment{Overlap threshold $\theta_o$ scales with $|S|$}
\State \Return \Call{FilterOverlap}{$\mathcal{D}$, $D_\text{prev}$, $\theta_\text{o}$, $\ell$}
\EndProcedure

\Require{Set of sets of candidate papers $\mathcal{D}$, set of previous papers $D_\text{prev}$, (inclusive) overlap threshold $\theta_o$, limit $\ell$}
\Ensure{Set of filtered papers $D_\text{out}$}
\Procedure{FilterOverlap}{$\mathcal{D}$, $D_\text{prev}$, $\theta_\text{o}$, $\ell$}
    \State $\mathcal{B} = \left[ \left(o, \{d \mid \left( \sum_{D \in \mathcal{D}} \mathds{1} (d \in D) \right) = o \}\right) \: \textbf{for} \: o \in \left[ |\mathcal{D}|, \dots, 1 \right] \right]$
    \Comment{Bin papers according to decreasing overlap}
    \State $D_\text{out} \leftarrow \{\}$
    \For{$o, D \in \mathcal{B}$}
    \Comment{Highest overlapping citations come first}
        \If{ $o < \theta_o \vee | D_\text{out} | \geq \ell$ }
            \textbf{break}
        \EndIf
        \State $D \leftarrow \{d \mid d \in D \wedge d \notin D_\text{prev} \}$
        \Comment{Filter out already present papers}
        \If {$\ell - |D_\text{out}| < |D|$}
        \Comment{If the entire bin won't fit within limit $\ell$}
            \State $D \leftarrow \{d \mid i, d \in $ \Call{$\text{sort}_{\downarrow \text{citers}}$}{$D$} $\wedge \: i \leq ( \ell - | D_\text{out} | ) \}$
            \Comment{Keep subset with the most future citers}
        \EndIf
        \State $D_\text{out} \leftarrow D_\text{out} \cup D$
    \EndFor
    \State \Return $D_\text{out}$
\EndProcedure



\end{algorithmic}
\end{algorithm}

\subsection{LitQA}

\subsubsection{Question Construction and Human Evaluation} \label{sec:question_evaluators}

LitQA questions were generated manually by a combination of the authors as well as contracted human experts (see \cite{laurent2024lab}). All human annotators were compensated, informed that their evaluations were being used in research of human-level performance, and consented to the use of their annotations and participation. Question authors were instructed to identify recent papers (published within the last ~36 months), and develop a multiple-choice question that requires context within the main text of the paper to answer and \textit{is not answerable by the abstract or title alone.} They were further advised that the question should require some amount of reasoning within the paper context, and not be a direct quote or statement from the paper. Distractors were instructed to be reasonable within the context, using either other information in the paper (e.g. other genes being discussed) or based on inherent or other knowledge. We also periodically tested question drafts with ChatGPT 3.5 or 4 (with logging disabled) to ensure questions were not easily answerable by models already, or to help design effective distractors by asking to provide plausible answers. Question drafts were also often searched against Google Scholar to ensure that it was not trivial to find the exact statement necessary to answer the question. This was an effective aid in revising question wording.

For contracted question authors or anyone new to drafting, the first 10 or so questions produced were carefully reviewed by one or more of the authors for quality, after which they were either asked to rework them, they were added to corpus, or they were removed from consideration. This feedback (referred to as \verb|calibration|) was usually enough to ensure quality question generation going forward from contractors, though questions were reviewed by authors on an ongoing basis prior to merging into the main corpus. Contractors were compensated via multiple structures during different phases of generation as we iterated on effective strategies. Initially, they were paid on an hourly basis at \$50 per hour. At later phases, they were paid on a completion basis such that they were paid per accepted question, which equated to variable hourly equivalents always \$50 or more per hour. Using this methodology, LitQA2 was built up from LitQA (47 questions) in two stages of releases, first 100 questions (147), then an additional 101 questions, adding to the original subset to make 248 total questions. 

Human evaluators were assigned questions from the LitQAv2 corpus in rounds of "quizzes" composed of 20-40 questions each. They were allocated up to a week to complete the quiz, but were given no other time constraints. They were also allowed to use tools such as internet search or journal collection search provided via their institutions. They were asked to explicitly refrain from using AI-based tools such as ChatGPT or Claude, though we did not have any method of enforcement of this request. In order to encourage completion as well as high performance, we attempted to design an incentive scheme that separately promoted both: In brief, we paid a base dollar amount (from \$3-6) to each question for completion. We then added a bonus to each question that was based on overall performance: 
\begin{itemize}
    \item{80\% or more overall score = bonus equal to base question amount for each correct question.}
    \item{60\% - 80\% overall score = bonus equal to half the base question amount for each correct question.}
    \item{Less than 60\% overall score = bonus equal to \$1 for each correct question.}
\end{itemize}

We additionally paid a completion bonus for the entire quiz in some instances, amounting to \$150. On average, evaluators were paid between \$50 to \$100 per hour based on self-estimation of time spent and total compensation. Quizzes were sometimes combined with similar evaluation quizzes for other evaluation benchmark categories we are developing.

For the human LitQA2 performance reported in this work, 2 rounds of quizzes (20 questions each) were given to 9 evaluators. In total, evaluators answered provided 266 unique answers across 248 unique questions.

A third round of quizzes was given with another set of 160 questions given to evaluators, with questions which overlapped from both the initial set of 147 questions and the remaining 101 questions. However, this quiz round had several outlier scores ( >90\%) in quizzes on the same subsets of questions which had been given previously to a much lower average (66.3\%). Upon investigating, we found that a portion of the initial 147 LitQA2 questions that were available on GitHub had been parsed by google-indexed data aggregators, causing it to trivially return results with a Google search from evaluators. Interestingly, this did not impact PaperQA2 performance, because PaperQA2 does not use Google websearch. Thus, we chose to exclude evaluator answers to the initial 147 LitQA2 questions in the third set of quizzes as it was no longer a valid human comparison.

\subsubsection{PaperQA measurement}
When measuring LitQA2 performance, \texttt{paperqa} function calls were performed for each LitQA question. LitQA answer order was randomized for each call, and, for all configurations shown, 3 full runs of all LitQA questions were performed. 3 runs were necessary to control for the inherent noise in LLM inference, even at 0 temperature (where all tests were performed).  LitQA2 was automatically evaluated using an evaluation LLM call (GPT-4-0613), which extracted the letter answer from PaperQA2's output. It was prompted to extract as follows: 

\begin{quote}
Extract the single letter answer from the following question and answer
\{prompt\_output(`QA')\}

Single Letter Answer:
\end{quote}

The extraction was then parsed and graded against the LitQA2 benchmarks correct ideal answer. All questions were graded as ``Unsure'' if the ``Insufficient information to answer this question'' option was selected, except for the LitQA2 questions where the ideal answer is ``null'', then the ``Insufficient information'' was graded as ``Correct''. 

LitQA2 was developed in two stages of 100 then 101 questions, adding to the original subset of LitQA (47) questions to make 248 questions. This gave us an opportunity to evaluate our system before and after adding the set of 101 previously unseen questions to evaluate overfitting. We can see that our system performed similarly in terms of accuracy between question sets, implying we were not overfitting to the original LitQA data. 

\begin{table}[h]
\centering
\begin{tabular}{|c|c|c|c|c|}
\hline
Ablation & Accuracy (Old 147) & Accuracy (New 101) & Precision (Old 147) & Precision (New 101) \\ \hline
Claude-3-Opus & 59.0\% $\pm $2.4\% & 54.5\% $\pm $6.1\% & 93.9\% $\pm $3.7\% & 82.6\% $\pm $6.0\% \\ \hline
Gemini-1.5-Pro & 55.8\% $\pm $5.7\% & 59.4\% $\pm $2.3\% & 93.2\% $\pm $1.8\% & 83.0\% $\pm $4.6\% \\ \hline
GPT-4-Turbo &  64.4\% $\pm $1.8\% & 63.7\% $\pm $7.4\% & 91.3\% $\pm $2.1\% & 81.8\% $\pm $7.3\% \\ \hline
\end{tabular}
\caption{Table of model choices using both stages of LitQA2 development, all errors shown are CI intervals.}
\label{tab:litqa_before_after}
\end{table}


\begin{figure}[h!]
    \centering % Center the figure
    \refstepcounter{suppfigure}
    \includegraphics[width=0.7\textwidth]{Figures/Wikicrow-Ablations.png}
    \caption{PaperQA2 precision performance on LitQA2 across all configuration categories included in \autoref{fig:litqa_performance}. All error bars are 95\% CI intervals.} % Add a caption
    \label{fig:litqa_precision} % Add a label for referencing the figure
\end{figure}

\begin{figure}[h!]
    \centering % Center the figure
    \refstepcounter{suppfigure}
    \includegraphics[width=0.9\textwidth]{Figures/Wikicrow-LitQAtest.png}
    \caption{Precision, accuracy and recall for different chunk size and parsing algorithm choices on LitQA. Note that the gap between recall and precision is because the model can sometimes find an alternative source to rule out distractors, provide the same information, or make the model confident enough to guess.} % Add a caption
    \label{fig:chunksize_ablations} % Add a label for referencing the figure
\end{figure}

When comparing PaperQA2 to other RAG technologies in \textbf{Figure  \ref{fig:litqa_performance}A}, the author's manually entered each LitQA2 question (in the exact format given to PaperQA2) into the online interfaces for both Perplexity\cite{perplexity} and Elicit\cite{elicit}. Perplexity was run using the "Pro" product, and GPT-4o, across all 248 LitQA2 questions. We recorded either the letter answer, or manually interpreted the response to assign the system's letter output. The performance metrics were calculated in an identical way to PaperQA. Elicit was not able to be run for all 248 questions, the numbers reported only represent the first tranche of 147 LitQA2 questions. On Elicit's final 101 LitQA2 questions, the authors noted a significant algorithm change. Answers resulted in > 90 \% "unsure" answers, thus we are only reporting on the subset of LitQA2 where multiple choice answers could be extracted. We also compared PaperQA2 to PaperQA's prior published configuration\cite{lala2023paperqa}, which used a smaller RCS model (GPT 3.5 Turbo) and ranking depth, had no citation traversal tool, and did not have access to Grobid parsings. PaperQA's LitQA2 accuracy was 36.7\% and its precision was 76.5\%. Thus, PaperQA2 showed a large improvement in both metrics, particularly in accuracy.

\subsection{WikiCrow} \label{sec:si-wikicrow}

WikiCrow statements were generated using a linked set of 5 different queries, 4 to PaperQA2, and 1 to a frontier LLM model, GPT-4-Turbo. As detailed in \textbf{Figure \ref{fig:WikiCrow_results}A}, each PaperQA2 query created a different article section. The PaperQA2 prompts and settings used for these queries can be found in our \texttt{paperqa} \href{https://github.com/Future-House/paper-qa}{repository}. The primary PaperQA2 parameter differentiation between WikiCrow and LitQA2 runs was the usage of Grobid for document parsing, as well as the modification of the RCS step to extract a gene name that is the focus of each paper chunk being summarized. They were designed to reduce hallucinations and accurately read data from tables.

We compared WikiCrow and Wikipedia statements starting from a sample of 300 genes which were sampled from from a complete list of all protein encoding human genes \cite{enwiki:proteinlist} (n=19,255) after being filtered for only those having valid (non-stub) Wikipedia articles (n=3,639). Among the 300 articles, Wikipedia statements were sampled, first from the pool of 300 genes, then among each article's \texttt{<p>} HTML sections with more than 25 characters. 500 total Wikipedia statements were sampled, which resulted in 240 unique genes. Matching WikiCrow articles were generated for each of the 240 genes for comparison. WikiCrow statements were broken at either a reference list (notated with parenthesis and PaperQA2 document names) or at a paragraph break, most closely matching Wikipedia's HTML structure. Each of the 500 Wikipedia statements were randomly matched with a WikiCrow statement from the same gene article to ensure that evaluators were given a similar gene-distribution of statements between WikiCrow and Wikipedia articles. 

References were extracted from the Wikipedia statements using the \texttt{<sup>} link, and matched to the correct DOI or hyperlink in the article references. References were extracted from the WikiCrow statements by parsing the parenthetical insertions (via regular expressions) and matching the DOI to the document name given by the model. In this way, each statement was associated with a list of DOI links to the sources listed. Efforts were then made to structure each statement similarly, removing any HTML from the Wikipedia statements, and removing all inserted references into the WikiCrow statements. All references were replaced with random numbers between 1 and 30, but kept consistent if listed several times in the same statement. References were then uniformly reformatted as \texttt{[x]}, to hide their origin.

The matching 500 Wikipedia and 500 WikiCrow statements were shuffled for each of the 4 evaluators. Each was given a list of 200 to grade while blinded to the statement origins. Graders were sent the questions with the following instructions: 

\begin{quote}
\textbf{Task overview}\\ \\
You’ll be reviewing statements written about human genes. These statements are intended to provide information taken from the literature, and should thus provide an accurate citation such that the stated information matches information contained in the cited literature. You’ll be evaluating each statement according to its accuracy as well as being cited appropriately. Please read the guidelines below carefully to ensure accurate work.\\ \\
\textbf{Scoring} \\

You’ll soon receive a link to a private Google Sheet with 200 statements, along with some associated information for each one, including links to any cited literature. You’ll be grading the sentences in each statement according to the below metrics. Some statements may contain multiple sentences, each of which should be evaluated such that any sentence not satisfying our scoring criteria counts against the entire statement. \\
\begin{itemize}
\item{Is it cited? In other words, does the statement contain one or more citations to published literature?}
\item{Is the information correct, as cited? In other words, is the information stated in the sentence correct according to the literature that it cites?\\}
\end{itemize}

There will be a column for each of the above scores. Put TRUE in the appropriate column if true according to the metric, and FALSE if not. There will also be a column for “notes”. Please use this column to include anything you think is important for your grading, or other flags we might want to know about. Below are some more detailed guidelines for scoring:

You are only evaluating if the statements are correct as cited -- not if the statement is globally correct or supported by other sources. Some common scenarios you may encounter:

\begin{itemize}

\item{\textbf{Simple statements:} Statement has a single citation and the content is supported by the citation. Is Cited = TRUE, Is Info Correct As Cited = TRUE}

\item{\textbf{Simple unsupported statements:} A statement is true, highly-specific (see "broad context" below, but not explicitly stated in the source, mark as FALSE for “correct as cited”, but mark "is cited" as TRUE.}

\item{\textbf{Unrelated, meta statements:} Statement is a meta-statement or unrelated to biology. For example, the statement may be the start of a list like: "Gene XYZ has the following properties:", or an attribution like "This data has been provided by Y organization." These statements should be judged as "Not Applicable/NA". They do not need to be cited.}

\item{\textbf{Broad context:} Statement has a citation and the content is a broader context than the citation, but not explicitly explained in the citation. For example, an explanation of what a secondary structure is or what a protein sequence is. The broad content should be judged to be "undergraduate biology student common knowledge". If it is, then the info can be marked correct as cited, if it's more advanced than undergraduate bio knowledge (or outright incorrect), then the info is incorrect as cited.}

\item{\textbf{Underspecified citation:} The statement may be factually correct, but the citation is under-specified. An example might be that the citation is a link to a database, where a user needs to query the database and synthesize their own data to support the statement. If it requires synthesis by the user, than the statement should be judged as NOT correct as cited. If it's a database link to the a page about the gene, and the fact can be found there, then it can be judges as correct as cited. }

\item{\textbf{Inaccessible citation:} The statement may be factually correct, but the citation is to an inaccessible source, like a closed source database or a textbook. These should be judged as "Not Applicable/NA". Though an attempt should be made to acquire the source to validate the information.}

\item{\textbf{Multi-part citations:} The statement has several sentences and citations. Each sentence should be correct as cited given all of the previous criteria (including the "Broad contexts" category).}

\item{Occasionally, it may not be clear from context which sentence applies to a source, there may be three sentences in a row with two sources trailing the sentences. In this case if the sentences are supported by either source, then it's correct as cited. If a single sentence is un-supported by either source, and they don't meet our broad context criteria, then mark as incorrect as cited.}

\item{If there are multiple citations, where one is inaccessible, but the accessible citations can support the statements by themselves, then mark as correct as cited. If the statement can NOT be validated by the accessible sources, then mark as "Not Applicable/NA".}
\end{itemize}
\end{quote}

Evaluators were incentivized to encourage both quality of review and quantity of statements reviewed. They were offered \$10 per statement graded as a base payment, and were additionally offered a \$10 per question bonus payment if they reviewed more than 50 questions by a specific deadline approximately five days after starting.

Using this criteria, 4 expert researchers were contracted as evaluators to grade 375 evaluations. They were instructed to only complete their statements in the order provided, and any statements that were not graded in a contiguous block from the initial statement were excluded from analysis (to avoid biasing towards simple statements). Among our 375 evaluations, 40 statements overlapped between 2 or more evaluators. Among those 40, 77.5\% (31) had agreement, and 22.5\% (9) had a disagreement in evaluation, showing good overall alignment in scoring. 

We had previously reported\cite{coxWikiCrow} lower ``cited and unsupported'' percentages for both Wikipedia and WikiCrow, but the prior study differed in key ways: 1. this work limits the statement samples to be from the same subset of gene articles, eliminating selection bias for longer, higher quality articles; 2. this work uses statements built from sentences which may cite one or more sources as is typically present in articles, the prior work focused on only single sentences with one source which greatly limited the sample; and 3. this work used a much more robust evaluation pipeline, with blinded external evaluators.

Due to the complexity and time required to evaluate WikiCrow accuracy, we relied on heuristics to evaluate our summarization-centric features, adding metadata (gene name) extraction to our RCS step and Grobid-based parsings. These features were meant to mitigate gene name confusion between the RCS step and the generate answer tool, and misreading tabular data, which anecdotally were the two most frequent issues found by the authors while developing WikiCrow. Among the reasoning issues reported by the evaluators, only 2 / 171 statements had issues with gene name confusion, and there were no reported table extraction errors. Utilizing Grobid structured parsings also reduced PaperQA2's token usage. Across a sample of 5,363 papers, PyMuPDF's parsings resulted in an average of 16,040 tokens per paper, while Grobid's parsings resulted in an average 8,903 tokens. The savings can be attribution to a reduction in excess whitespace as well as the removal of reference sections, which we did not include in our Grobid parsings.

\subsection{Contradiction Detection Methods}  \label{sec:cc-methods}
ContraCrow detects contradictions in literature through two steps. Unless provided, the first step in ContraCrow is to extract claims from a given paper \textbf{\autoref{fig:contradiction_detection}A}. First, the provided paper is split into chunks. To maintain context for extracting claims, each chunk is only split within its respective section, up to 5000 characters (no overlap), and the section and paper titles are included. Each resulting chunk is then fed into a claim-extraction LLM to extract claims candidates. The resulting extracted claims are each evaluated by a filtering LLM that assesses the claims based on quality and generalizability, assigning a score out of 10. Only claims scoring 8 or higher are considered. 

Once claims have been generated or are provided, each claim is fed into PaperQA2 with a special \textit{contradiction detection prompt}, instructing the system to search for contradictions and respond with an appropriate format. The system considers each claim independently and identifies relevant papers across literature in order to identify possible contradictions to the claim. Finally, as instructed by the \textit{contradiction detection prompt}, ContraCrow outputs its reasoning and a choice from an 11-point Likert scale (see \autoref{fig:contradiction_detection}D). We map the natural language score to integers (0-10) in our calculations, which allows us to tune the decision threshold, as seen in \ref{fig:contradiction_detection}B. For all experiments, this step uses Claude 3.5 Sonnet \cite{claude2024} as the LLM model, a chunk size of 7000 characters, a temperature of 0, and a simple overlap of 250. Unless otherwise specified, we use a decision boundary of $8$ on this step. 

\subsubsection{ContraDetect} \label{sec:cc-contradetect}
The ContraDetect benchmark was generated from the LitQA2 contradiction dataset, removing any questions with ``null'' responses as to not overlap with the no-evidence data. The resulting questions were randomly split into two groups. The questions and corresponding ideal answers from the first group were individually given to \texttt{gpt-4-turbo-2024-09}\cite{achiam2023gpt} with instructions to rephrase them into factual statements. The second group was similarly turned into incorrect statements (contradictions), given question and the first corresponding distractor answers. Each question corresponds to exactly 1 statement. These claims then proceeded into ContraCrow's second step, as described above.

We separately designed 42 no-evidence statements by hand. The authors carefully designed questions in their respective fields of expertise and did extensive literature review to ensure that no literature has reported on the claim. There is no overlap between the no-evidence statements and the LitQA2 questions. The resulting answers from ContraCrow were also evaluated by the authors to ensure that no evidence had been found. Any claims with found evidences were removed from this dataset. 

\subsubsection{Human Evaluation} \label{sec:cc-eval}
We randomly selected 100 papers from our local database of biology papers. 7 of these papers failed to parse, leaving our dataset at 93 papers. These papers have no overlap with the papers corresponding to LitQA2 questions. \texttt{gpt-4-turbo-2024-04-09}\cite{achiam2023gpt} was used for both LLM models in the claim-generation step. The resulting claims were then fed into the contradiction-detection step as described above, in batches of 1,000 claims.

For the human ``contradiction validation'' task, 50 claims scored at 8 and 50 claims at least 9 by ContraCrow were randomly chosen (distribution: $8:50$, $9:48$, $10:2$). The resulting 100 contradictions were split over 5 expert evaluators evenly, with no overlap. The evaluators were provided with the claim, the original claim information (source paper reference and the chunk used to generate the claim), some background information (acronyms, definitions, etc.), ContraCrow's reasoning, and all chunks cited in the reasoning. The evaluators were tasked with giving two separate binary responses indicating their agreement with both the model's reasoning and conclusion, based on the provided cited chunks (up to 10). They were instructed to use additional provided resources sparingly, only for clarification, and to \textit{only} consider the evidence chunks and ContraCrow's reasoning when making a final decision. Only the label for ``agreement with the model's conclusion'' was used for analysis. 

For the human ``contradiction detection'' task, 10 claims scored 0-4, 10 claims scored 8, and 10 claims scored at least 9 by ContraCrow were randomly chosen (distribution: $0:1$, $1:6$, $2:2$, $3:1$, $8:10$, $9:10$). The resulting 30 claims were split over 5 evaluators. Each claim was evaluated by at least 2 evaluators. The evaluators were provided with the claim, the original claim information (source paper reference and the chunk used to generate the claim), background information (acronyms, definitions, etc.), and all literature chunks considered by ContraCrow (up to 15). The evaluators were tasked with assigning a Likert scale value and a binary (``yes'' or ``no'') label to determine whether each claim contained a contradiction to the provided sources. For this task, the evaluators did not have access to ContraCrow's reasoning. Similar to the first task, they were instructed to use any provided materials for clarification, but to only use the provided literature when making a final decision. 

In the both the ``contradiction validation'' task and the binary ``contradiction detection'' task, ContraCrow scores of $>=8$ were considered the positive case in order to compare to humans' binary outputs. For the Likert score tasks, raw scores were used.

Evaluators were compensated for their work on grading claims to encourage thorough review. For the ``contradiction validation'' task, they were paid \$12 per claim completed. For ``contradiction validation'', they were paid \$20 per claim completed. They were additionally paid a \$200 `completion bonus' for completing grading on all claims by a deadline approximately one week after starting.

Evaluators were sent the questions with the following instructions:  \\
\begin{quote}
\textbf{Contradiction Detection Instructions} \\
\textbf{Task overview}\\ 
For this task, we are interested in identifying and rating contradicting research findings, or claims. We’ve assembled a set of claims from research papers that are potentially contradicted by other findings, either prior or later. Your task will be to both assess the claim itself, as well as review portions of text taken from research papers that may be related to and may or may not present contradicting information to the presented claim.
We will not have an explicit calibration phase for this task, nor do we provide examples. More detailed instructions are provided below to help guide you. You may use any sources available to you in order to understand concepts, acronyms, etc. that you are unfamiliar with, but you should perform the evaluation itself by-hand and using only the claim and context provided.
We also explicitly ask that you do not use ChatGPT or other AI tools to come to your conclusions, though they may be used if helpful to define terms or clarify concepts. Please do not put claims or chunks into ChatGPT and similar tools.

\textbf{Instructions}\\
This task will consist of two separate activities, Claim Evaluation and Contradiction Detection (type a and b). For each claim, you will be given the following information:
\begin{itemize}
    \item Claim: a single sentence making some claim or finding.
    \item Chunk: the section of text from the source paper that was used to generate the claim.
    \item Section: the section of the paper that the chunk lives in, if available
    \item Title: the title of the paper that the chunk lives in, if available
    \item DOI: the DOI URL for the paper that the chunk lives in, if available
\end{itemize}

\textbf{Section 1: Claim Evaluation}\\
Your first task is to a determination of the claim being a valid claim. Some examples of good \& bad claims are shown below:
\begin{itemize}
    \item Good claims are typically clear, testable, and generalizable
    \begin{itemize}
        \item ``Cells expressing ARG1 are more likely to undergo apoptosis.''
    \end{itemize}
    \item Bad claims will come in different forms, like not having enough context, i.e. it does not make sense outside the context of the paper.
    \begin{itemize}
        \item ``Cells in cluster 1 had higher expression of ARG1.''
        \item ``Cells expressing this gene are more likely to undergo apoptosis.''
        \item ``Method A performed better than Method B in our study.''
        \item Claims provided may also just not really be defined as a `claim'. These might describe methods, descriptions, or obvious common knowledge
            \begin{itemize}
                \item ``We used method A to test the effects of ARG1 on cell death.''
                \item ``ARG1 is a protein.''
            \end{itemize}
        \item Claims may also not be supported by the provided chunk (Mixed up gene names, off-topic discussion, etc.)
    \end{itemize}
\end{itemize}
To evaluate each claim, select all that apply from the following options:
\begin{itemize}
    \item good claim
    \item not enough context
    \item not generalizable
    \item not a claim
    \item not supported
\end{itemize}

Some important considerations:
\begin{itemize}
    \item Multiple claims may come from the same chunk \& claims may be redundant in nature
    \item Consider each claim independently
    \item You may use the all provided information in this section to help you on this step (including the paper title, as it may provide context)
    \item Avoid:
    \begin{itemize}
        \item Evaluating the claim for truth value
        \item Evaluating the wording/syntax of the claim, unless it changes the meaning
        \item Deep-diving into the paper outside of the claim
        \item Spending more than a few minutes grading each claim, on average
    \end{itemize}
\end{itemize}

\textbf{Section 2: Contradiction Detection}\\
This part will contain additional information for each claim, consisting of up to 15 ``evidence chunks'' which are sections of research paper text that may serve as evidence either supporting, contradicting or neutral to the claim to various degrees.
Each chunk will itself contain (when available) the source paper's citation. This information can be used for reference and clarity (for defining acronyms, clarifying things, etc) but should not be considered when determining contradictions.
Your job for this part of the task is to determine whether (and to what extent) the claim contradicts or agrees with the provided evidence chunks. You will provide a selection for both of the below categories (note there are two separate selections!)
\begin{itemize}
    \item Contradiction grade (choose one of the following):
        \begin{itemize}
            \item Explicit Agreement
            \item Strong Agreement
            \item Agreement
            \item Possibly an Agreement
            \item Lack of Evidence
            \item Possibly a Contradiction
            \item Nuanced Contradiction
            \item Contradiction
            \item Strong Contradiction
            \item Explicit Contradiction
        \end{itemize}
    \item Contradiction determination:
        \begin{itemize}
            \item YES - Claim contradicts the provided context
            \item NO - Claim does not contradict the provided context
        \end{itemize}
\end{itemize}
Some important considerations:
\begin{itemize}
    \item A claim may have multiple chunks from the same paper.
    \item Consider claims individually. If the claim is a bad claim, use your best judgement about whether or not and how to evaluate the contradiction. We ask that you attempt to evaluate it if possible, and if you feel you cannot proceed, please explain why then move on to the next claim.
\item Focus on the claim + context provided: Don't worry about other contradictions or agreements in the context, only focus on the specific claim. Additional information (publication information, background info) is provided for clarification where needed, but you should rely solely on the provided context and claim for labeling.
\item It is possible that the provided statement makes more than 1 claim. Use your best judgement in this case, considering the context provided.
\item Avoid:
    \begin{itemize}
        \item Evaluating the claim or evidences for truth value or quality (only evaluate in relation to each other).
        \item Deep-diving into any text outside of the provided context (including the papers from which the context was derived).
    \end{itemize}
\end{itemize}
\end{quote}

\begin{quote}
\textbf{Contradiction Validation Instructions }\\
\textbf{Task overview}\\
This task is very similar, but this time you will be given an AI model's reasoning on the evidence chunks, and only the evidence chunks the model considered (which should cut down the review time considerably). You will do the usual annotation, and additionally say whether you agree with the LLM's reasoning. Section 1 is unchanged, Section 2 is different.

We also explicitly ask that you do not use ChatGPT or other AI tools to come to your conclusions, though they may be used if helpful to define terms or clarify concepts. Please do not put claims or chunks into ChatGPT and similar tools.
\\
\textbf{Instructions}\\
This task will consist of two separate activities, Claim Evaluation and Contradiction Detection (type a and b). For each claim, you will be given the following information:
\begin{itemize}
    \item Claim: a single sentence making some claim or finding.
    \item Chunk: the section of text from the source paper that was used to generate the claim.
    \item Section: the section of the paper that the chunk lives in, if available
    \item Title: the title of the paper that the chunk lives in, if available
    \item DOI: the DOI URL for the paper that the chunk lives in, if available
\end{itemize}

\textbf{Section 1: Claim Evaluation}\\
\textit{\*\*Section 1 is the same as in the ``contradiction detection'' task.\*\*}

\textbf{Section 2: Contradiction Detection - Model Evaluation}\\
This part will contain additional information for each claim. Each claim will have a model response, which is some reasoning to answer the question `Are there any contradictions to this claim in literature?' The model response will reference up to 10 “evidence chunks” which are sections of research paper text.
These referenced chunks will also be provided for you and can be linked via the citation key in the model response (for example, the model response might reference \texttt{Song2022DEPTH2:pages 16-17}. The chunk corresponding to this reference will also be found under \texttt{Song2022DEPTH2:pages 16-17}.
Each chunk will itself contain (when available) the source paper's citation. This information can be used for reference and clarity (for defining acronyms, clarifying things, etc) but should not be considered when evaluating contradictions or model response.
Your job for this task is to determine whether you agree with the model response and to what extent the claim contradicts or agrees with the evidence chunks provided.
You will provide a selection for both of the below categories (note there are two separate selections!)
\begin{itemize}
    \item Contradiction grade (choose one of the following):
        \begin{itemize}
            \item Explicit Agreement
            \item Strong Agreement
            \item Agreement
            \item Possibly an Agreement
            \item Lack of Evidence
            \item Possibly a Contradiction
            \item Nuanced Contradiction
            \item Contradiction
            \item Strong Contradiction
            \item Explicit Contradiction
        \end{itemize}
    \item Contradiction determination:
        \begin{itemize}
            \item YES - Claim contradicts the provided context
            \item NO - Claim does not contradict the provided context
        \end{itemize}
    \item Optional Explanation: optionally, you may provide a written explanation/justification of your scoring.
\end{itemize}

Some important considerations:
\begin{itemize}
    \item A claim may have multiple chunks from the same paper.
    \item Consider claims individually.
    \item Because the model response is provided, we ask that you evaluate every provided claim.
    \item Additional information (publication information, background info) is provided for clarification where needed.
    \item It is possible that the provided statement makes more than one claim. Use your best judgment in this case, considering the model response provided.
    \item Avoid:
    \begin{itemize}
        \item Evaluating the claim or evidences for truth value or quality (only evaluate in relation to each other).
        \item Evaluating the model response truth value, quality, syntax, or style, except where directly related to the claim.
        \item Deep-diving into any text outside of the provided context (including the papers from which the context was derived).
    \end{itemize}
\end{itemize}
\end{quote}

\bibliographystyle{unsrt}  
\bibliography{references}

\end{document}


\clearpage
\appendix

\section{Appendix A: Task Descriptions}
\label{app:tasks}

These cards give enough context to interpret the results; the builders that assemble each family's tasks and splits are in the released code.

\begin{taskbox}{MBPP+\cardsub{validation task $\cdot$ anchor: hidden unit tests}}
\small
\textbf{Task.} A one-paragraph natural-language programming problem plus one example call; the solver writes a Python function.\\[2pt]
\textbf{Anchor.} The EvalPlus hidden test suite, which adds many more tests per problem than the original MBPP and repairs faulty ones, giving reliable pass/fail.\\[2pt]
\textbf{Typical failures.} Crashes, mishandled edge cases, printing instead of returning, misread specifications, exactly the failure classes the op pool checks.
\end{taskbox}

\begin{taskbox}{Spider 2.0-Snow\cardsub{validation task $\cdot$ anchor: execution match}}
\small
\textbf{Task.} An analyst question over a real hosted Snowflake warehouse: production-scale databases with hundreds of tables, cryptic columns, sharded table families, and dialect-specific SQL. The solver writes one query from the question, a rendered schema catalog, and optional domain notes.\\[2pt]
\textbf{Anchor.} Execute the candidate and compare the fetched table to the stored gold result with the official comparator (value-based; column names and order are ignored).\\[2pt]
\textbf{Typical failures.} Unresolvable identifiers, type errors on date-like columns, and, hardest, queries that run cleanly but compute the wrong thing. Our single-call solver passes about a third of screened tasks.
\end{taskbox}

\begin{taskbox}{Report generation\cardsub{target task $\cdot$ no accurate metric exists}}
\small
\textbf{Task.} Write one \emph{section} of an analyst report (e.g., ``value-chain map: control points from ASIC to optical engine'') in markdown, given the inputs below.\\[2pt]
\textbf{Inputs.} (i)~Section title and purpose. (ii)~A style contract fixing units, reference dates, and accounting rules (e.g., ``all amounts USD, calibrated to 2026~Q2, market sizes must not be summed across pools''). (iii)~Evidence cards. (iv)~A metric vocabulary registering the section's key figures (id, name, value, unit).\\[2pt]
\textbf{Evidence card.} An audited unit of evidence: id \texttt{E0103}, one-sentence claim, verbatim source quote, source attribution and date, stance (supporting / opposing / neutral), and an auditor's credibility verdict.\\[2pt]
\textbf{Requirements.} Ground every claim in the cards (cite \texttt{[E\#\#\#]}); state each key figure's value with its inline tag binding it to the metric vocabulary (rendered downstream as audit tooltips); hedge weak or disputed evidence and present both sides of conflicts; include required structure such as a summary line.\\[2pt]
\textbf{Why no metric exists.} Section quality is holistic; even an expert author cannot reduce ``a good section'' to a checkable rule. The RAQS rubric scores lexical and structural proxies of the requirements above and is treated as reference signal only.
\end{taskbox}

\section{Appendix B: Setup Details}
\label{app:setup}

\paragraph{Splits and screening.}
For MBPP+ and Spider~2.0 we screen every task with five independent samples of the frozen single-call solver, graded by the official anchor. We drop tasks the solver always passes (no room to improve) and stratify the rest across pass levels for a balanced difficulty mix, then split 60/40 into train and test. The dev set is ten items, five clear failures and five clean passes (for report generation: five golden-quality positives and five class-stratified checkable failures), so the metric is scored on both classes. In the released code the dev and test splits are named \texttt{eval\_dev} and \texttt{eval\_locked}.

\paragraph{Anchors.}
\emph{MBPP+}: hidden unit tests. \emph{Spider~2.0}: execute the candidate on the live warehouse and compare the fetched table to the stored gold result with the official comparator, which matches values (ignoring column names and order). We reproduce the evaluator's CSV round-trip of the fetched frame, without which even gold queries fail on dtype. \emph{Report generation}: the metric loop's anchor is a frozen demo-delta soft label, where a teacher compares the frozen-agent report against the golden section under a closed seven-class miss taxonomy and passes iff at most one miss class is present. Each of the 48 held-out sections contributes two locked items, the frozen-agent report (which the strict anchor almost always fails) and the golden section (which passes by definition), giving 96 items balanced 50/50, so a degenerate metric scores 0.5. The skill and co-loops are evaluated by the official rubric evaluator (RAQS, 0--100 normalized), whose dimensions cover evidence grounding, metric-tag discipline, claim calibration, structure, and style.

\paragraph{Spider~2.0 prompt grounding.}
The generation prompt grounds the solver with a compact typed schema catalog (columns as name:TYPE, sharded tables collapsed, sample values on date-like columns) and Snowflake-specific rules (quoted identifiers, type-respecting casts, value-only grading). We report this because it is a large, shared effect we control for, not a contribution: grounding raises the frozen baseline by half on the screened pool (0.211 to 0.329) but is applied identically to every arm, so all reported lifts sit on top of it. It also sets up the detectability contrast: grounding pushes the solver's residual failures from compile errors into clean-executing wrong values.

\section{Appendix C: The Metric Language and Lifecycle}
\label{app:lang}

\paragraph{Language.} Recapping the grammar of Eq.~\ref{eq:grammar}: a metric is an expression over ops (pure functions of (task, output) returning \textsc{drawback}/\textsc{clean}/\textsc{abstain}), composed by $\bigvee$, $\bigwedge$, negation, the unweighted $K$-of-$k$ vote, and a weighted vote against a threshold, with abstaining children excluded from a combinator and an all-abstain node abstaining. Negation and the weighted vote are offered to the composer on every round but appear in none of the selected metrics, so every reported expression reads as a plain drawback disjunction or vote.

\paragraph{Lifecycle.} Each task starts from a small pool of seed detectors, the obvious failure checks a practitioner reads straight off the task description (does the code crash, does the query group correctly, is a key figure missing); drafting them, or having an LLM draft them from the task spec, is an afternoon's work, and their exact contents barely matter because the lifecycle grows and prunes the pool from there. Synthesis authors typed op specs from clustered misses under a closed per-task taxonomy. The birth gate requires a new op to fire on at least half its cluster and stay clean on known-good outputs. Ops born from anchored misses enter active; unlabeled-gap ops enter shadow and are promoted only on demonstrated anchor agreement. Per-op fitness is the leave-one-out marginal of the incumbent expression, and non-positive ops retire after a two-round grace.

\paragraph{Selection and guards.} The objective multiplies dev agreement by consensus agreement on unlabeled train outputs (exponent 1.0 for MBPP+ and Spider, 0.25 for report generation, where op consensus is a weaker regularizer), minus a small complexity penalty. Both agreement terms are weighted. Dev agreement is the recall-weighted mean of the two per-class recalls over the dev items a candidate opines on, with the \textsc{fail} class at weight 2 and the \textsc{clean} class at weight 1. Train consensus is reliability-weighted: each opining op's vote is scaled by $1 + 10 \max(0, m_o)$ for that op's anchored leave-one-out marginal $m_o$, so anchor-proven ops lead the vote and ops with no anchored evidence contribute at unit weight. This keeps a bloc of cheaply synthesized, correlated detectors from defining the pseudo-label that the regularizer rewards agreeing with. The complexity penalty charges 0.001 per op beyond a floor of three and a further 0.002 per LLM-calling op, which is why selected expressions stay small and mostly deterministic. Fail-closed anchoring makes candidates without a usable dev opinion unselectable, and a validity gate drops all-pass/all-fail/all-abstain candidates. The naive ablation disables exactly these two guards; the no-lifecycle ablation disables exactly the birth gate, shadow tier, and retirement.

\section{Appendix D: Final-Judge Protocol and Rubrics}
\label{app:judge}

\paragraph{Protocol.} For each of the 48 held-out sections we pair the stored round-0 report against the stored peak-round report, and separately against the final-round report, for every arm and seed. Pre-repair this is 3 arms $\times$ 3 seeds $\times$ 2 rounds $\times$ 48 sections $=$ 864 pairs; the post-repair rerun adds a fourth seed on the two co-loop arms, giving 1{,}056. Each pair is judged twice with positions swapped, and a win counts only when both orders agree; otherwise it is a tie.

\paragraph{Rubrics.} The generic rubric asks a senior-editor question about holistic reader value with no knowledge of the pipeline's conventions. The task-aware rubric additionally states the format contract: evidence citations and inline metric tags are required output that renders downstream; a tag with no value nearby is a defect; fabricated figures are disqualifying. Both rubrics judge the same stored pairs, so the four cells of Table~\ref{tab:judge} differ only in the repair pass (rows) and the rubric (columns).

\section{Appendix E: What Evolution Actually Produced}
\label{app:artifacts}

Final artifacts from the paper runs, verbatim except where marked as condensed, illustrating the central qualitative claim: everything the loops evolve is legible. Figure~\ref{fig:evolved} gives the final metric on each task, Figure~\ref{fig:oplife} traces one synthesized detector through the whole lifecycle, Figure~\ref{fig:grading} grades one attempt with the hidden tests and with the evolved metric side by side, and Figure~\ref{fig:samelesson} pairs skills the two loops learned independently.

\begin{figure*}[t]
\begin{taskbox}{Evolved metrics\cardsub{one final expression per task{,} verbatim}}
\small
\setlength{\tabcolsep}{0pt}
\renewcommand{\arraystretch}{1.15}
\begin{tabular}{@{}>{\raggedright\arraybackslash}p{0.30\linewidth}@{\hspace{14pt}}>{\raggedright\arraybackslash}p{0.29\linewidth}@{\hspace{14pt}}>{\raggedright\arraybackslash}p{0.34\linewidth}@{}}
\textbf{MBPP+} & \textbf{Spider 2.0} & \textbf{Report} \\
\addlinespace[2pt]
Flag a solution if the judge finds a spec mismatch, or it crashes on probe inputs, or it prints instead of returning:
&
Flag a query if it aggregates without grouping, or the judge finds it does not answer the question (selected independently on two of three seeds):
&
Flag a section if a registered key figure never appears, or a weak-credibility card is cited without hedging (the other two seeds swap the first leaf for an uncited-claims detector):
\\
\addlinespace[4pt]
\texttt{(any}\par
\texttt{~~spec\_mismatch}\par
\texttt{~~crash}\par
\texttt{~~returns\_not\_print\_only)}
&
\texttt{(any}\par
\texttt{~~missing\_group\_by}\par
\texttt{~~spec\_sql\_mismatch)}
&
\texttt{(any}\par
\texttt{~~ledger\_metric\_uncovered}\par
\texttt{~~low\_status\_card\_overclaimed)}
\\
\end{tabular}
\end{taskbox}
\caption{\textbf{The final evolved metric on each task, verbatim, with a plain reading above it.} Each is a short drawback disjunction over typed detectors: nothing here is a learned scalar, and each leaf names the failure class it fires on.}
\label{fig:evolved}
\end{figure*}

\begin{figure*}[t]
\begin{taskbox}{A synthesized op\cardsub{born{,} gated{,} promoted{,} and selected (Spider 2.0)}}
\small\raggedright
Tracing one op through the whole lifecycle, on the metric-loop seed that ends at \texttt{(any syn\_forbids\_select\_star\_when\_columns spec\_sql\_mismatch)}. The loop clustered four unlabeled train outputs that returned every column although the question implied specific ones, and authored the typed spec \texttt{\{kind: static, check: forbids\_select\_star, detects: select\_star\_specific\_columns\}}. Because its birth evidence was an unlabeled gap rather than an anchored miss, it entered as a \emph{shadow}: it ran and recorded verdicts but could not be selected. It was promoted the same round only once composing it with the incumbent raised dev agreement from 0.71 to 0.81, and selection then adopted it, replacing the grouping check the seed pool had contributed. The two co-loop seeds on this task converge on a different synthesized op (\texttt{syn\_top\_n\_missing\_limit}) and the third on the seed \texttt{select\_star} check, so the leaf is not a fixed artifact of the pool.
\end{taskbox}
\caption{\textbf{One synthesized detector through the whole lifecycle: born from unlabeled outputs, admitted as a shadow, promoted only once it raised anchor agreement, then selected.} A Spider~2.0 metric-loop seed; the aggregate counts are Table~\ref{tab:oplife}.}
\label{fig:oplife}
\end{figure*}

\begin{figure*}[t]
\begin{taskbox}{Grading one task: hidden tests vs.\ evolved metric\cardsub{MBPP+ task{,} final-round trace{,} verbatim}}
\small
\setlength{\tabcolsep}{0pt}
\renewcommand{\arraystretch}{1.15}
\begin{tabular}{@{}>{\raggedright\arraybackslash}p{0.36\linewidth}@{\hspace{14pt}}>{\raggedright\arraybackslash}p{0.27\linewidth}@{\hspace{14pt}}>{\raggedright\arraybackslash}p{0.31\linewidth}@{}}
\textbf{Task and solver output} & \textbf{Ground-truth anchor} & \textbf{Evolved metric} \\
\addlinespace[2pt]
``Write a function to find the ratio of zeroes in an array of integers.''
\par\vspace{4pt}
\texttt{def zero\_count(nums):}\par
\texttt{~~zeros = sum(1 for x}\par
\texttt{~~~~~~~~~~in nums if x == 0)}\par
\texttt{~~non\_zeros = len(nums) - zeros}\par
\texttt{~~...}\par
\texttt{~~return zeros / non\_zeros}
\par\vspace{4pt}
The code runs cleanly and returns a number; it computes the ratio to \emph{non-zeros} where the spec wants the ratio to \emph{all} elements. Only the meaning is wrong.
&
Hidden unit tests (EvalPlus).
\par\vspace{4pt}
\texttt{assert zero\_count(...)}\par
\texttt{~~$\rightarrow$ AssertionError}
\par\vspace{4pt}
Verdict: \textsc{fail}.
\par\vspace{4pt}
The verdict is reliable but opaque: no indication of which requirement was misread.
&
Evolved composition \texttt{(any spec\_mismatch crash returns\_not\_print\_only)}.
\par\vspace{4pt}
\texttt{crash}: clean\par
\texttt{returns\_not\_print\_only}: clean\par
\texttt{spec\_mismatch} (judge op): \textsc{drawback}
\par\vspace{4pt}
Verdict: \textsc{fail}, agreeing with the anchor, and \emph{naming} the failure class: the specification was misread.
\\
\end{tabular}
\par\vspace{5pt}
The named drawback is what turns this attempt into a usable failure capsule for skill synthesis; the anchor's bare verdict could not have driven the same repair.
\end{taskbox}
\caption{\textbf{The same verdict, reached two ways: the hidden tests say \textsc{fail}, and the evolved metric says \textsc{fail} while naming why.} One MBPP+ attempt whose code runs cleanly and returns a number but computes the wrong quantity. Only the middle column is ground truth; the right column is what an inspectable composition adds on top of an equally correct verdict.}
\label{fig:grading}
\end{figure*}

\begin{figure*}[t]
\begin{taskbox}{Same lesson{,} two loops\cardsub{functionally matched skills{,} reference vs.\ co-evolved{,} condensed}}
\small
Each pair below was synthesized independently, one by the reference loop under ground-truth grading and one under the evolved metric (Double Ratchet); neither loop saw the other. Converging on the same task lesson is the qualitative face of the retention result.
\par\vspace{5pt}
\setlength{\tabcolsep}{0pt}
\renewcommand{\arraystretch}{1.15}
\begin{tabular}{@{}>{\raggedright\arraybackslash}p{0.31\linewidth}@{\hspace{14pt}}>{\raggedright\arraybackslash}p{0.31\linewidth}@{\hspace{14pt}}>{\raggedright\arraybackslash}p{0.31\linewidth}@{}}
\textbf{MBPP+} & \textbf{Spider 2.0} & \textbf{Report} \\
\emph{(grader imports a literal identifier)} & \emph{(empty result diagnosis)} & \emph{(the anti-erasure lesson)} \\
\addlinespace[4pt]
\emph{reference} \ \texttt{function-name-} \texttt{from-assert}: ``the harness imports the symbol literally as written in the assert; lift the identifier verbatim including case, the prose is only a hint.''
&
\emph{reference} \ \texttt{sql-empty-result-} \texttt{diagnosis}: ``a 0-row result is almost never correct when the gold has rows; incrementally relax predicates to locate which condition eliminates all rows.''
&
\emph{reference} \ \texttt{metric-tag-annotates-} \texttt{not-replaces}: ``the scorer counts a figure only when a literal number sits adjacent to the brace tag; a tag with no rendered numeral reads as an unrendered placeholder.''
\\
\addlinespace[6pt]
\emph{co} \ \texttt{copy\_callable\_} \texttt{identifier\_from\_assert}: ``the grader imports the name verbatim from the assert; reproduce that exact spelling letter for letter, including apparent typos.''
&
\emph{co} \ \texttt{sql-empty-result-} \texttt{sanity-probe}: ``an empty result almost always means one predicate matched nothing; validate each CTE's row count in isolation before trusting the join.''
&
\emph{co} \ \texttt{metric-tag-inline-} \texttt{after-number}: ``the tag annotates the value; it never replaces the value, and a bare card cite does not satisfy it.''
\\
\end{tabular}
\end{taskbox}
\caption{\textbf{Trained on ground truth and trained on the evolved metric, the two loops independently learn the same task lesson.} One functionally matched skill pair per task, condensed. Neither loop saw the other, so the convergence is the qualitative face of the lift-retention result in Table~\ref{tab:main}.}
\label{fig:samelesson}
\end{figure*}

\section{Appendix F: Lifecycle Statistics}
\label{app:lifecycle}

Table~\ref{tab:oplife} in the results gives the metric-loop op lifecycle; Table~\ref{tab:skilllife} here gives the skill-loop side, both for the anchored arms. The locked metric sets are small, 40 items on MBPP+ and Spider and 96 on reports, but they are regraded in full every round: the nine anchored metric-loop runs record 52{,}600 item-level verdicts against the locked reference and 5{,}743 birth-gate decisions. Each run stores every verdict with its per-op trace, which is what makes the paired tests recomputable from a finished run without rerunning any loop; the released code includes the script that reads those tables and computes the reported counts and $p$-values. For the skill loops, synthesis is prolific (roughly 1.5 skills per round), the bounded bank forces contribution-based eviction of more than half of everything created, and rollback fires rarely; the report skill loop is the outlier, its targeted failure capsules producing skills so effective that a handful suffice (one seed held its plateau with a bank of six).

\begin{table}[t]
\centering
\footnotesize
\setlength{\tabcolsep}{2.2pt}
\begin{tabular}{llcccc}
\toprule
Task & Loop & Synthesized & Bank & Evicted & Rollback \\
\midrule
MBPP+ & skill & 140--167 & 50 & 80--100 & 1 \\
 & \textbf{co} & \textbf{111--154} & \textbf{50} & \textbf{60--101} & \textbf{0} \\
Spider & skill & 149--177 & 50 & 86--114 & 1--2 \\
 & \textbf{co} & \textbf{152--153} & \textbf{50} & \textbf{76--86} & \textbf{1--2} \\
Report & skill & 6--66 & 6--50 & 0--15 & 0 \\
 & \textbf{co} & \textbf{135--156} & \textbf{50} & \textbf{85--106} & \textbf{0} \\
\bottomrule
\end{tabular}
\caption{\textbf{Synthesis is prolific and the bounded bank evicts more than half of what it creates.} Skill-loop lifecycle over 100 rounds, bank cap 50, the counterpart to Table~\ref{tab:oplife}. \emph{Bank}: final active skills. \emph{Evicted}: removed by contribution. \emph{Rollback}: held-out regressions reverted. Three-seed min--max; a bare number means all seeds agree. Bold rows are our proposed configuration.}
\label{tab:skilllife}
\end{table}

\section{Appendix G: Negative and Supporting Results}
\label{app:negative}

\paragraph{Naive co-loops are not a free pass.}
Although naive co-loops match the anchored arm on held-out evaluation (their measurement anchor is untouched), their metric arms are degenerate where inspection matters: on report generation the naive metric grades 0.97--1.00 of train as passing (versus 0.78--0.82 anchored), and on Spider~2.0 the naive objective barely moves on two of three seeds.

\paragraph{Skill-loop robustness to metric quality.}
The weak-metric Spider~2.0 seeds reported in the results reach the same peak as their two-op sibling because training signal quality degrades gracefully: a failed attempt still carries concrete error text into skill synthesis even when the pass/fail label is noisy.

\paragraph{Anchor size.}
Ten dev items is a design choice about what sparse anchoring realistically supplies, not a tuned setting, so we test how far selection depends on it. The results section reports the numbers; the protocol behind them is an offline replay of report-generation selection, drawing 200 stratified subsamples per anchor size from the same ten items, preserving the fail/clean balance within each draw, and rescoring every candidate expression the live run actually considered under the reduced anchor. No new generation and no new solver calls are involved, so the comparison isolates the anchor's contribution to selection from run-to-run variance.

\section{Appendix H: Reproducibility}
\label{app:repro}

The code is released at \url{https://github.com/amazon-science/Self-Evolving-Agents-Double-Ratchet}: all loops, op packs, split builders, screening scripts, and the judge protocols, sufficient to rerun every experiment end to end and reproduce results in line with those reported here. We release code only, not run data, results, or logs. All model calls go through the Amazon Bedrock API, so wall-clock time is dominated by API latency. Each run writes its own round-by-round history, op pool, skill bank, and judge outputs as structured result files, and every number in this paper is read directly from those files. The role-to-model table named in the setup covers every LLM role: solver, teacher, composer, skill synthesizer, judge ops, and final judge. Spider~2.0 requires credentialed access to the benchmark's hosted warehouse; all other experiments are self-contained.

\end{document}
\typeout{get arXiv to do 4 passes: Label(s) may have changed. Rerun}
